\chapter{The Network-on-Chip}\label{ch:noc}
\begingroup\emergencystretch=3em

\whatyoufind{The flit-level network that connects L1s, LLC banks and memory controllers on a
grid, a torus, a 3-D mesh or torus, a ring, a star or any graph you write down. What it models, from packets and flits
to routers, credits and virtual channels. The routing algorithms and arbiters it ships. How the
directory protocols stay free of deadlock on it. LLC banks and memory controllers as network
endpoints. How to configure and run it, what it reports, and how it is validated.}

\section{What the network models}\label{sec:noc-overview}

The network is the interconnect \code{Network}. The system \code{MultiCoreSystem\_Mesh} selects it
when \key{interconnect\_type} is \code{Network}. It replaces the single shared medium of a bus with routers
joined by point-to-point links (\cref{fig:noc-system}). Every cache, LLC bank and memory
controller is an \emph{endpoint}. Each endpoint has a \emph{network interface} (NI) that attaches
it to one router through a local port. The controllers keep the \code{CommunicationInterface}
contract of every other interconnect. They push and pop messages as before and never learn that a
network sits between them.

A message travels as a \emph{packet}. A packet is cut into \emph{flits}, the unit a link carries
in one tick and a buffer slot holds. The model is flit-level wormhole switching with virtual
channels (VCs) and credit-based flow control. A head flit reserves a path through each router,
the body and tail flits follow it, and a flit moves only when the buffer ahead of it has room.
The blocking that dominates worst-case latency in a real network is therefore simulated, not
estimated.

\begin{figure}[htbp]\centering
\begin{adjustbox}{max width=\textwidth}% System view of the flit-level network: the 2x2 preset MultiCoreSystem_Grid2B2M (L1s, two LLC
% banks with their memory ports, two memory controllers) and how a 4x4 grid generalises it.
% Port numbers are the ones Topology.cpp assigns: link ports east, west, north, south (absent
% ones skipped), then the local ports in ascending endpoint id.
\begin{tikzpicture}[font=\sffamily\scriptsize, >=Latex, x=1mm, y=1mm,
  rtr/.style={draw=octpurple, fill=octio, very thick, minimum width=13mm, minimum height=13mm, inner sep=0, align=center, font=\sffamily\scriptsize\bfseries},
  ni/.style={draw=octink, fill=octquiet!25, minimum width=4mm, minimum height=3.6mm, inner sep=0, font=\sffamily\tiny\bfseries},
  ep/.style={draw=octrule, fill=octmem, rounded corners=1.5pt, minimum width=13mm, minimum height=4.4mm, inner sep=0.5pt, font=\sffamily\scriptsize\bfseries},
  bank/.style={ep, fill=octctl, draw=octamber},
  mport/.style={ep, fill=octctl, draw=octamber, densely dashed},
  mc/.style={ep, fill=white, draw=octteal, thick},
  wire/.style={octink, semithick, ->},
  loc/.style={octink, thin},
  pnum/.style={font=\sffamily\tiny\bfseries, text=octpurple, inner sep=0.6pt},
  ttl/.style={font=\sffamily\small\bfseries, text=octink},
  sub/.style={font=\sffamily\scriptsize, text=octquiet, align=left},
  srtr/.style={draw=octpurple, fill=octio, thick, minimum size=6mm, inner sep=0, font=\sffamily\scriptsize\bfseries}]
  % ================= left: the 2x2 grid =================
  \node[ttl, anchor=west] at (-27, 17) {2$\times$2 \code{grid2d}, preset \code{MultiCoreSystem\_Grid2B2M}};
  \node[rtr] (R0) at (0,0)    {R0\\[-0.3ex]\mdseries (0,0)};
  \node[rtr] (R1) at (40,0)   {R1\\[-0.3ex]\mdseries (1,0)};
  \node[rtr] (R2) at (0,-38)  {R2\\[-0.3ex]\mdseries (0,1)};
  \node[rtr] (R3) at (40,-38) {R3\\[-0.3ex]\mdseries (1,1)};
  % links: one per direction
  \draw[wire] ($(R0.east)+(0,2)$) -- ($(R1.west)+(0,2)$);
  \draw[wire] ($(R1.west)+(0,-2)$) -- ($(R0.east)+(0,-2)$);
  \draw[wire] ($(R2.east)+(0,2)$) -- ($(R3.west)+(0,2)$);
  \draw[wire] ($(R3.west)+(0,-2)$) -- ($(R2.east)+(0,-2)$);
  \draw[wire] ($(R0.south)+(-2,0)$) -- ($(R2.north)+(-2,0)$);
  \draw[wire] ($(R2.north)+(2,0)$) -- ($(R0.south)+(2,0)$);
  \draw[wire] ($(R1.south)+(-2,0)$) -- ($(R3.north)+(-2,0)$);
  \draw[wire] ($(R3.north)+(2,0)$) -- ($(R1.south)+(2,0)$);
  \node[text=octquiet, align=center] at (20,-19) {two links per\\connection,\\one per direction};
  % link port numbers at each router
  \node[pnum, anchor=south west] at ($(R0.east)+(0.3,3)$) {0 E};
  \node[pnum, anchor=south east] at ($(R1.west)+(-0.3,3)$) {W 0};
  \node[pnum, anchor=north west] at ($(R0.south)+(3,-0.3)$) {1 S};
  \node[pnum, anchor=south west] at ($(R2.north)+(3,0.3)$) {1 N};
  \node[pnum, anchor=north west] at ($(R1.south)+(3,-0.3)$) {1 S};
  \node[pnum, anchor=south west] at ($(R3.north)+(3,0.3)$) {1 N};
  \node[pnum, anchor=south west] at ($(R2.east)+(0.3,3)$) {0 E};
  \node[pnum, anchor=south east] at ($(R3.west)+(-0.3,3)$) {W 0};
  % endpoints of R0 (west side): L1 0, bank 10, memory port 110 -> local ports 2, 3, 4
  \foreach \e/\st/\yy/\pp in {L1 0/ep/4/2, LLC 10/bank/-0.5/3, port 110/mport/-5/4} {
    \node[ni] at (-11,\yy) {NI};
    \node[\st, anchor=east] at (-14,\yy) {\e};
    \draw[loc] (-9,\yy) -- (-6.5,\yy) node[pnum, pos=0.5, above] {\pp};
  }
  % endpoints of R2 (west side): L1 2, memory controller 101 -> 2, 3
  \foreach \e/\st/\yy/\pp in {L1 2/ep/-35.5/2, MC 101/mc/-40.5/3} {
    \node[ni] at (-11,\yy) {NI};
    \node[\st, anchor=east] at (-14,\yy) {\e};
    \draw[loc] (-9,\yy) -- (-6.5,\yy) node[pnum, pos=0.5, above] {\pp};
  }
  % endpoints of R1 (east side): L1 1, memory controller 100 -> 2, 3
  \foreach \e/\st/\yy/\pp in {L1 1/ep/2.5/2, MC 100/mc/-2.5/3} {
    \node[ni] at (51,\yy) {NI};
    \node[\st, anchor=west] at (54,\yy) {\e};
    \draw[loc] (46.5,\yy) -- (49,\yy) node[pnum, pos=0.5, above] {\pp};
  }
  % endpoints of R3 (east side): L1 3, bank 11, memory port 111 -> 2, 3, 4
  \foreach \e/\st/\yy/\pp in {L1 3/ep/-34/2, LLC 11/bank/-38.5/3, port 111/mport/-43/4} {
    \node[ni] at (51,\yy) {NI};
    \node[\st, anchor=west] at (54,\yy) {\e};
    \draw[loc] (46.5,\yy) -- (49,\yy) node[pnum, pos=0.5, above] {\pp};
  }
  % legend
  \begin{scope}[shift={(-27,-54)}]
    \node[ep, minimum width=9mm, anchor=west] at (0,0) {L1};     \node[sub, anchor=west] at (10,0) {cache};
    \node[bank, minimum width=9mm, anchor=west] at (30,0) {LLC}; \node[sub, anchor=west] at (40,0) {LLC bank, a home};
    \node[mport, minimum width=9mm, anchor=west] at (0,-6) {port}; \node[sub, anchor=west] at (10,-6) {an LLC bank's memory port};
    \node[mc, minimum width=9mm, anchor=west] at (48,-6) {MC};   \node[sub, anchor=west] at (58,-6) {memory controller};
    \node[ni, anchor=west] at (2.5,-12) {NI}; \node[sub, anchor=west] at (10,-12) {network interface, one per endpoint};
    \node[rtr, minimum width=5mm, minimum height=4mm, anchor=west] at (60,-12) {}; \node[sub, anchor=west] at (66,-12) {router};
    \node[pnum, anchor=west] at (3.5,-18) {3};
    \node[sub, anchor=north west, text width=72mm] at (10,-15.6) {port number at that router: link ports first (E, W, N, S, absent ones skipped), then one local port per endpoint in ascending id};
  \end{scope}
  % ================= right: 4x4 generalisation =================
  \begin{scope}[shift={(92,4)}]
    \node[ttl, anchor=west] at (-9, 13) {4$\times$4 \code{grid2d}};
    \foreach \yy in {0,1,2,3} \foreach \xx in {0,1,2,3} {
      \pgfmathtruncatemacro{\rr}{\yy*4+\xx}
      \node[srtr] (s\rr) at (\xx*12, -\yy*12) {\rr};
    }
    \foreach \yy in {0,1,2,3} \foreach \xx in {0,1,2} {
      \pgfmathtruncatemacro{\aa}{\yy*4+\xx}\pgfmathtruncatemacro{\bb}{\yy*4+\xx+1}
      \draw[octink, semithick] (s\aa) -- (s\bb);
    }
    \foreach \yy in {0,1,2} \foreach \xx in {0,1,2,3} {
      \pgfmathtruncatemacro{\aa}{\yy*4+\xx}\pgfmathtruncatemacro{\bb}{\yy*4+\xx+4}
      \draw[octink, semithick] (s\aa) -- (s\bb);
    }
    % axes
    \draw[->, octquiet] (-6,6) -- (41,6) node[right] {$x$};
    \draw[->, octquiet] (-6,6) -- (-6,-41) node[below] {$y$};
    \foreach \xx in {0,1,2,3} \node[text=octquiet] at (\xx*12, 8.3) {\xx};
    \foreach \yy in {0,1,2,3} \node[text=octquiet] at (-8.3, -\yy*12) {\yy};
    % port numbers of two routers
    \node[pnum, anchor=west] at ($(s5.east)+(0.2,1.6)$) {0};
    \node[pnum, anchor=east] at ($(s5.west)+(-0.2,1.6)$) {1};
    \node[pnum, anchor=south] at ($(s5.north)+(1.8,0.2)$) {2};
    \node[pnum, anchor=north] at ($(s5.south)+(1.8,-0.2)$) {3};
    \node[pnum, anchor=west] at ($(s0.east)+(0.2,1.6)$) {0};
    \node[pnum, anchor=north] at ($(s0.south)+(1.8,-0.2)$) {1};
    \node[sub, anchor=north west, text width=60mm] at (-9,-46) {routers row-major: $r = y\cdot k_x + x$\\[0.3ex]
      corner (0): E\,0, S\,1, local from 2\\
      edge (1): E\,0, W\,1, S\,2, local from 3\\
      interior (5): E\,0, W\,1, N\,2, S\,3, local from 4\\[0.3ex]
      keys: \code{dims(vi),4,4}, one\\\code{node\_map[<id>](i)} per endpoint};
  \end{scope}
\end{tikzpicture}
\end{adjustbox}
\caption{The network of \code{MultiCoreSystem\_Grid2B2M}, the largest system preset: four L1s,
two LLC banks with their memory ports and two memory controllers on a 2$\times$2 grid. Each
endpoint has its own NI and local port. Right: the same numbering on a 4$\times$4 grid.}
\label{fig:noc-system}
\end{figure}

\begin{table}[htbp]\centering\small
\caption{Terms used in this chapter.}\label{tab:noc-terms}
\begin{tabularx}{\textwidth}{@{}lY@{}}
\toprule term & meaning \\ \midrule
endpoint & a cache, an LLC bank, an LLC bank's memory port or a memory controller, with its own id \\
NI & network interface: the endpoint's side of the network, its injection and RX queues \\
router & a switching node with link ports and local ports \\
port & one side of a link at a router. Link ports are numbered first, local ports last \\
link & a one-way channel between two routers. A connection is two links, one per direction \\
packet & one message for one destination, as a sequence of flits \\
flit & what a link carries in one tick and a buffer slot holds, \key{flit\_bytes} wide \\
class & REQUEST, FORWARD or RESPONSE: a virtual network with its own buffers \\
VC & virtual channel: one flit FIFO at an input port, belonging to one class \\
credit & permission to send one flit into one downstream VC \\
tick & one step of the network. One tick is one core cycle as shipped \\
home & the LLC bank that owns a line's directory entry \\
owner & the endpoint a message is attributed to (\code{Message::owner}) \\
$R$, $f$ & routers a packet crosses (first and last included), flits in a packet \\
\bottomrule
\end{tabularx}
\end{table}

\paragraph{Scope.} The network carries the directory protocols: directory MSI (the presets),
MESI and MOESI. Their traffic is point-to-point, and three message classes keep it free of
deadlock (\cref{sec:noc-deadlock}). Snooping protocols do not run on the network, because they need
an ordered broadcast that a mesh does not give. Real-time arbitration on the network (TDM slot
tables, priority VCs, a \code{TDM\_XY} routing) is not available. Its key names are reserved and
refused at start-up. \Cref{sec:noc-limits} lists every limitation.

\section{Running it}\label{sec:noc-running}

Six presets ship (\cref{tab:noc-presets}). The system \code{MultiCoreSystem\_Mesh} runs every
one of them except the network-only one, with the preset given as its configuration file:

\begin{lstlisting}[style=shell]
# directory MSI on the 2x2 grid, one LLC on router 3
./build/Octopus_Simulator -s MultiCoreSystem_Mesh -c MultiCoreSystem_Grid \
    -p "workload_path(s)=BMs/eembc-traces/a2time01-trace/" \
    -o results/grid
# the same with the LLC moved to router 0 and round-robin allocation
./build/Octopus_Simulator -s MultiCoreSystem_Mesh -c MultiCoreSystem_Grid \
    -p "interconnect.node_map[10](i)=0" \
    -p "interconnect.arbiter_type(s)=RRArbiter" -o results/grid-llc0
# two banks and two memory controllers, each simulated by MCsim
./build/Octopus_Simulator -s MultiCoreSystem_Mesh -c MultiCoreSystem_Grid2B2M \
    -p "main_memory_type(s)=MCsim" -o results/grid2b2m
# the network alone: sixteen traffic generators on a 4x4 grid
./build/Octopus_Simulator -s NetworkOnly -o results/netonly
\end{lstlisting}

Every network key lives under \code{interconnect.} in the system file and can be overridden with
\code{-p} like any other key (\cref{ch:configuration}). Their defaults are in
\file{configuration/Interconnect/Network.csv}, which extends \file{Mesh.csv}.

\begin{table}[htbp]\centering\small
\caption{The network presets in \code{configuration/SystemConfigurations/}. The four
\code{Grid} presets are 2$\times$2 with \code{XY} routing, three classes, injection and RX queues
of 32, and \code{requires\_ordering pair\_per\_class}. \code{NetworkOnly} is 4$\times$4 with the
default queues of 16 and no ordering requirement. \code{NetNoC} is a \code{custom} graph with an
injection queue of 1024.}\label{tab:noc-presets}
\begin{tabularx}{\textwidth}{@{}lY@{}}
\toprule preset & what it contains \\ \midrule
\code{MultiCoreSystem\_Grid} & L1 0 to 3 on routers 0 to 3, one LLC (id 10) on router 3 beside L1 3. Directory MSI. Memory on a point-to-point bus off the grid \\
\code{MultiCoreSystem\_Grid2B} & two LLC banks, 10 on router 0 and 11 on router 3, lines interleaved, each bank half the single LLC. Both banks share one memory bus \\
\code{MultiCoreSystem\_Grid2BM} & \code{Grid2B} with the memory controller on the network: memory ports 110 and 111 beside the banks, controller 100 on router 1 \\
\code{MultiCoreSystem\_Grid2B2M} & \code{Grid2BM} with a second controller, 101 on router 2, and 4\,KiB pages interleaved across the two \\
\code{MultiCoreSystem\_NetNoC} & the legacy single-switch NoC expressed as a network (\cref{sec:noc-legacy}) \\
\code{NetworkOnly} & a 4$\times$4 grid with a traffic generator on every router and no caches (\cref{sec:noc-generators}) \\
\bottomrule
\end{tabularx}
\end{table}

To run directory MESI or MOESI on a grid preset, set four protocol keys as on the bus directory
system (\cref{ch:coherence-engine}). For MESI they are:

\begin{lstlisting}[style=shell]
-p "cache_controller[*].protocol_type(s)=DIRECTORY_MESI" \
-p "cache_controller[*].fsm_filename(s)=MESI_directory_v2" \
-p "llc_controller.protocol_type(s)=DIRECTORY_LLC_MESI" \
-p "llc_controller.fsm_filename(s)=MESI_LLC_directory_v2"
\end{lstlisting}

MOESI uses \code{DIRECTORY\_MOESI}, \code{MOESI\_directory\_v2}, \code{DIRECTORY\_LLC\_MOESI} and
\code{MOESI\_LLC\_directory\_v2}. On the banked presets (\code{Grid2B}, \code{Grid2BM},
\code{Grid2B2M}) the two LLC keys take the prefix \code{llc\_controller[*].} instead of
\code{llc\_controller.}, so that every bank gets them.

\section{Topology and endpoint attachment}\label{sec:noc-topology}

A topology is a set of routers, the links between them and the router each endpoint attaches to.
\key{topology} names the \emph{builder} that makes it. Seven builders ship:

\begin{description}
\item[\code{grid2d}] $k_x \times k_y$ routers from \key{dims(vi),kx,ky}, numbered row-major,
  $r = y\cdot k_x + x$, with $y$ growing downwards (north is $y-1$)
  (\cref{fig:noc-system}, right).
\item[\code{torus2d}] the grid plus wrap links on every row and column. Both dimensions must be
  at least 3, so that a wrap link differs from a direct link.
\item[\code{mesh3d}] $k_x \times k_y \times k_z$ routers from \key{dims(vi),kx,ky,kz}, numbered
  $r = x + k_x\cdot(y + k_y\cdot z)$. Layer $z$ is a $k_x \times k_y$ grid numbered as
  \code{grid2d}.
\item[\code{torus3d}] the 3-D mesh plus wrap links in all three dimensions, each at least 3.
\item[\code{ring}] $k \ge 3$ routers from \key{dims(vi),k}, a one-dimensional torus.
\item[\code{star}] one router with every endpoint on it.
\item[\code{custom}] any graph from \key{connection\_matrix[i]} rows (node id, then its
  neighbours), as for the mesh. The matrix must be symmetric. A node in \key{component\_ids} is an
  endpoint and gets an \emph{implicit router} of its own, which stands for the endpoint's link
  ports and adds no router latency. The other nodes are explicit routers (\key{router\_ids} may
  list them).
\end{description}

The five generated shapes other than \code{star} are \emph{lattices}: each router has a coordinate
per dimension, a dimension may wrap, and each link port has a direction. Routing algorithms read
this capability rather than a topology name, so an algorithm written for lattices works on all of
them (\cref{sec:noc-routing}).

\paragraph{Port order.} A router numbers its link ports first, in a fixed order, and its local
ports last (\cref{tab:noc-ports}). On a lattice the order follows the directions, and a direction
that does not exist at the edge of a mesh is skipped, so the numbers shift. A wrap link takes the
place of the missing direction and is named with \code{(wrap)}, for example \code{east(wrap)}. On
\code{custom} a node's link ports follow the order in which its own row names its neighbours. In
\code{MultiCoreSystem\_NetNoC} the row \code{0,200,10} gives L1 0 port 0 to switch 200 and port 1
to the LLC.

\begin{table}[htbp]\centering\small
\caption{Link port order per builder. Local ports follow the link ports.}\label{tab:noc-ports}
\begin{tabularx}{\textwidth}{@{}lY@{}}
\toprule builder & link ports, in order \\ \midrule
\code{grid2d}, \code{torus2d} & east ($x+1$), west ($x-1$), north ($y-1$), south ($y+1$) \\
\code{mesh3d}, \code{torus3d} & east, west, north, south as on the grid, then up ($z-1$) and down ($z+1$) \\
\code{ring} & clockwise ($r+1$), counter-clockwise ($r-1$) \\
\code{custom} & the neighbours in the order of the node's own matrix row. A neighbour named twice gets one port \\
\code{star} & none \\
\bottomrule
\end{tabularx}
\end{table}

Port numbers name ports and carry no routing intent. A preferred route on a \code{custom} graph
is stated with \key{routing\_via} (\cref{sec:noc-routing}). \code{NetNoC} uses it to send L1-to-L1
traffic through switch 200 rather than leaving the choice to a tie-break by the lowest port.

\paragraph{Endpoints.} On the generated topologies \key{node\_map[<id>](i)} names the router of
every endpoint in \key{component\_ids}, and a missing entry stops the run with the key named.
Several endpoints may share a router. With \key{local\_ports} 0, the default, each endpoint gets
its own local port, numbered after the link ports. On the lattices the local ports follow
ascending endpoint id, on \code{star} and \code{custom} the order of \key{component\_ids}. In
\code{MultiCoreSystem\_Grid} router 3 has west (0), north (1), L1 3 on local port 2 and the LLC on
local port 3. A positive \key{local\_ports} smaller than the endpoints at a router makes them
share ports round-robin. The NIs on a shared port then compete for it through the port's arbiter,
and the RX credits stay per endpoint and class.

\section{Packets, flits and classes}\label{sec:noc-packets}

\begin{figure}[htbp]\centering
\begin{adjustbox}{max width=\textwidth}% From a message to flits, to a class (virtual network), to a virtual channel of the local input
% port. Numbers are the shipped defaults: header_bytes 8, flit_bytes 16, 64-byte lines.
\begin{tikzpicture}[font=\sffamily\scriptsize, >=Latex, x=1mm, y=1mm,
  hdr/.style={draw=octink, fill=octquiet!35, minimum height=5mm, inner sep=0, anchor=west},
  pay/.style={draw=octink, fill=octmem, minimum height=5mm, inner sep=0, anchor=west},
  flt/.style={draw=octink, fill=white, minimum height=5mm, minimum width=13.6mm, inner sep=0, anchor=west, font=\sffamily\scriptsize\bfseries},
  hflt/.style={flt, fill=octquiet!35},
  q/.style={draw=octink, fill=white, minimum height=4mm, minimum width=4mm, inner sep=0},
  box/.style={draw=octrule, rounded corners=2pt, fill=octpaper, inner sep=3pt, align=left},
  cls/.style={draw=octink, minimum width=20mm, minimum height=5mm, inner sep=1pt, font=\sffamily\scriptsize\bfseries},
  lbl/.style={text=octquiet},
  ttl/.style={font=\sffamily\small\bfseries, text=octink, anchor=west},
  arr/.style={->, thick, octink}]
  % ---------- row 1: bytes to flits ----------
  \node[ttl] at (0, 10) {1. A message becomes a packet of flits};
  \node[lbl, anchor=west] at (0, 3.5) {data message (a 64-byte line)};
  \node[hdr, minimum width=7.2mm] at (0,-2) {8};
  \node[pay, minimum width=57.6mm] at (7.2,-2) {payload \code{data\_size} = 64 bytes};
  \draw[decorate, decoration={brace, amplitude=3pt, mirror}] (0,-5) -- (64.8,-5) node[midway, below=2pt, lbl] {\code{header\_bytes} + payload = 72 bytes};
  \node[hflt] at (0, -18) {head};
  \foreach \i in {1,2,3} \node[flt] at (\i*14.4, -18) {body};
  \node[flt] at (57.6, -18) {tail};
  \draw[decorate, decoration={brace, amplitude=3pt, mirror}] (0,-21) -- (13.6,-21) node[midway, below=2pt, lbl, anchor=north west, xshift=-7mm] {\code{flit\_bytes} = 16};
  \node[lbl, anchor=west] at (72.5, -18) {$\lceil 72/16 \rceil = 5$ flits};
  \node[lbl, anchor=west] at (0, -32) {request, forward or acknowledgement (no data)};
  \node[hdr, minimum width=7.2mm] at (0,-38) {8};
  \draw[arr] (8.5,-38) -- (16.5,-38);
  \node[hflt, minimum width=19mm] at (17,-38) {head = tail};
  \node[lbl, anchor=west] at (38, -38) {$\lceil 8/16 \rceil = 1$ flit};
  % formula
  \node[box, anchor=north west, text width=63mm] at (99, 7) {\textbf{Flit count} (\code{flitCount} in \code{NetTypes.h})\\[0.6ex]
    $\textit{payload} = \textit{data\_size}$ if the message carries data, else 0\\[0.8ex]
    \hspace*{8mm}$f = \left\lceil \dfrac{\textit{header\_bytes} + \textit{payload}}{\textit{flit\_bytes}} \right\rceil$\\[0.8ex]
    With \code{flit\_bytes = 0} every packet is one flit: the message-level model.\\[0.6ex]
    The head flit carries the route. Body and tail flits follow it in order and are never interleaved with another packet on the same VC (wormhole).};
  % ---------- row 2: class and VC ----------
  \node[ttl] at (0, -47) {2. The packet joins the queue of its class, then a VC of that class};
  \node[box, anchor=north west, text width=33mm] (cl) at (0,-56) {\textbf{Classifier}\\by \code{Message::kind}, then fallback rules on \code{data}, \code{from}, \code{to} and the home ids};
  \node[cls, fill=white] (c0) at (56,-58) {0 REQUEST};
  \node[cls, fill=octquiet!20] (c1) at (56,-66) {1 FORWARD};
  \node[cls, fill=octquiet!45] (c2) at (56,-74) {2 RESPONSE};
  \foreach \c in {c0,c1,c2} \draw[arr] (cl.east) -- (\c.west);
  \node[lbl, anchor=south] at (82,-55.5) {NI injection queues};
  \foreach \yy in {-58, -66, -74} \foreach \k in {0,1,2,3} \node[q] at (76+\k*4.2, \yy) {};
  \foreach \c/\yy in {c0/-58, c1/-66, c2/-74} \draw[arr] (\c.east) -- (73.8, \yy);
  \node[lbl, anchor=north, align=center] at (82,-77) {\code{inject\_queue\_depth}\\packets per class};
  \draw[arr] (91,-66) -- (104,-66) node[midway, above, lbl] {1 flit};
  \node[lbl, anchor=south, align=center] at (132,-55.5) {local input port of the router};
  \foreach \c/\yy in {0/-57, 1/-65, 2/-73} {
    \foreach \v/\dy in {0/0, 1/-3.4} {
      \foreach \k in {0,1,2,3} \node[q, minimum height=3mm, minimum width=3.4mm] at (126+\k*3.6, \yy+\dy) {};
      \node[lbl, anchor=east, font=\sffamily\tiny] at (124, \yy+\dy) {class \c, vc \v};
    }
  }
  \node[lbl, anchor=north, align=center] at (132,-77) {\code{[class][vc]} FIFOs\\\code{vc\_depth} flits each};
  \node[box, anchor=north west, text width=158mm] at (0,-87) {Each tick the local port sends one flit, chosen by the port's arbiter among every class (oldest first under \code{FCFSArbiter}). A new packet takes the lowest-index free VC of its class that its first route admits (\code{vc\_mask}) and that holds a credit. The drawing shows \code{vcs\_per\_class} 2. The shipped default is 1.};
\end{tikzpicture}
\end{adjustbox}
\caption{A message becomes a packet of flits, the classifier gives it a class, and the NI queues
it for that class before it takes a VC of the class at the router's local input port.}
\label{fig:noc-packet}
\end{figure}

A message with several destinations (an LLC invalidating its sharers) becomes one packet per
destination, in the order of its \code{to} list. The NI accepts all of them or none. Each packet
carries the whole \code{Message}, which is delivered unchanged. The only field the network writes
is \code{from}, stamped at injection with the sending endpoint's id, as every Octopus interface
does. A packet's flit count is (\cref{fig:noc-packet})
\[
f = \left\lceil \frac{\textit{header\_bytes} + \textit{payload}}{\textit{flit\_bytes}} \right\rceil ,
\]
where the payload is the message's \code{data\_size} if it carries data and 0 otherwise. With the
defaults (\key{header\_bytes} 8, \key{flit\_bytes} 16) a request is one flit and a 64-byte line is
five. \key{flit\_bytes} 0 makes every packet one flit, the message-level model.

\paragraph{Classes.} Every packet belongs to one class. A class is a virtual network: it has its
own VCs at every input port, its own injection queue and its own RX queue at every NI. With
\key{num\_classes} 3, the default, the classifier reads \code{Message::kind}
(\cref{tab:noc-classes}). An untagged message falls back on its shape: with data it is a
RESPONSE, without data to a home it is a REQUEST, and without data from a home, or owned by a
home, it is a FORWARD. A message that matches no rule stops the run
with its fields printed. \key{num\_classes} 2 applies the legacy split of the mesh: class 0 for
messages without data and class 1 for messages with data, whatever their kind.

\begin{table}[htbp]\centering\small
\caption{Message kinds and their classes with three classes.}\label{tab:noc-classes}
\begin{tabularx}{\textwidth}{@{}llY@{}}
\toprule class & index & kinds \\ \midrule
REQUEST & 0 & GetS, GetM, PutS, PutM, PutE, PutO, any demand request \\
FORWARD & 1 & forwarded requests (FwdGetS, FwdGetM), invalidations, a home's acknowledgement without data (PutAck, the MOESI ack count), a bank's memory read (\code{MEM\_READ}) \\
RESPONSE & 2 & data responses and supplies, InvAck, fills, write-back data, memory writes (\code{MEM\_WRITE}), write acknowledgements (\code{MEM\_WACK}) \\
\bottomrule
\end{tabularx}
\end{table}

A home's PutAck travels in the FORWARD class on purpose. The directory L1 handles a home's Inv
followed by its PutAck to the same cache, but not the reverse order. In one class with ordered
routing (\cref{sec:noc-routing}) the PutAck can never overtake the Inv. An acknowledgement is
consumed without sending anything, so the FORWARD class stays drainable.

\section{The network interface}\label{sec:noc-ni}

\begin{figure}[htbp]\centering
\begin{adjustbox}{max width=\textwidth}% The network interface between a controller and its router: transmit-side admission (ADMIT-TX),
% injection, ejection with the RX reservation, RX serving, and the endpoint admission property.
\begin{tikzpicture}[font=\sffamily\scriptsize, >=Latex, x=1mm, y=1mm,
  ctl/.style={draw=octamber, fill=octctl, rounded corners=2pt, thick},
  nib/.style={draw=octink, fill=octquiet!12, rounded corners=2pt, thick},
  rtr/.style={draw=octpurple, fill=octio, rounded corners=2pt, very thick},
  stp/.style={draw=octrule, fill=white, rounded corners=1.5pt, text width=35mm, align=left, inner sep=2pt, anchor=north west},
  q/.style={draw=octink, fill=white, minimum height=3.2mm, minimum width=3.2mm, inner sep=0},
  qf/.style={q, fill=octquiet!45},
  qr/.style={q, fill=octamber!40},
  arr/.style={->, thick, octink},
  lbl/.style={text=octquiet},
  ttl/.style={font=\sffamily\small\bfseries, text=octink, anchor=west},
  box/.style={draw=octrule, rounded corners=2pt, fill=octpaper, inner sep=3pt, align=left}]
  % frames
  \draw[ctl] (0,6) rectangle (41,-68);  \node[font=\sffamily\small\bfseries, anchor=north west] at (1,5) {controller};
  \draw[nib] (47,6) rectangle (128,-68); \node[font=\sffamily\small\bfseries, anchor=north west] at (48,5) {network interface (NI)};
  \draw[rtr] (134,6) rectangle (165,-68); \node[font=\sffamily\small\bfseries, anchor=north west, text=octpurple] at (135,5) {router};
  % ---------------- transmit ----------------
  \node[stp] (a1) at (2,-1) {\textbf{ADMIT-TX}: pop a message only if the injection queues its transition will use have room (\code{txFreeSlots}), then reserve them};
  \node[stp] (a2) at (2,-24) {run the actions: each push uses a reserved slot and never fails};
  \draw[arr] (a1) -- (a2);
  \node[lbl, anchor=west] at (48,-1) {\textbf{transmit}};
  \node[stp, text width=27mm] (p1) at (49,-4.5) {\code{pushMessage}: stamp \code{from}, one packet per destination, class per packet. All or nothing.};
  \draw[arr] (a2.east) -- (44,-29) |- (p1.west);
  \node[lbl, anchor=west] at (81,-1) {per class, \code{inject\_queue\_depth}};
  \foreach \c/\yy in {REQ/-6, FWD/-11, RSP/-16} {
    \node[lbl, anchor=east, font=\sffamily\tiny, inner sep=0.5pt] at (93.2,\yy) {\c};
    \foreach \k in {0,1,2} \node[q] at (95+\k*3.4,\yy) {};
    \node[qr] at (105.2,\yy) {};
    \node[qf] at (108.6,\yy) {};
    \draw[arr] (110.4,\yy) -- (117.2,-11);
    \draw[arr] (p1.east |- 0,-11) -- (87.2,\yy);
  }
  \node[lbl, align=left, anchor=north west, text width=25mm] at (85,-19.5) {amber: reserved by an admitted transition\\grey: oldest packet};
  \node[draw=octink, fill=white, circle, inner sep=0.8pt, font=\sffamily\tiny\bfseries] (arb) at (119,-11) {arb};
  \node[lbl, align=center, anchor=north] at (119,-14) {oldest\\eligible\\packet};
  \draw[arr] (arb.east) -- node[above, lbl, pos=0.3] {1 flit} (134,-11);
  \node[lbl, align=left, anchor=north west, text width=28mm] at (135,-7) {local input port: \code{[class][vc]} FIFOs, \code{vc\_depth} credits each};
  % ---------------- receive ----------------
  \node[lbl, anchor=west] at (48,-37) {\textbf{receive}};
  \node[lbl, anchor=west] at (80,-40) {RX per class, \code{eject\_queue\_depth}};
  \foreach \c/\yy in {RSP/-44, FWD/-49, REQ/-54} {
    \node[lbl, anchor=west, font=\sffamily\tiny] at (107,\yy) {\c};
    \node[qf] at (89,\yy) {}; \foreach \k in {1,2} \node[q] at (89+\k*3.4,\yy) {};
    \node[qr] at (99.2,\yy) {}; \node[q] at (102.6,\yy) {};
  }
  \node[lbl, align=left, anchor=north west, text width=40mm] at (86,-57) {grey: next to serve\\amber: reserved for a packet in flight};
  \draw[arr] (134,-49) -- node[above, lbl, pos=0.62] {tail ejected} (113,-49);
  \node[lbl, align=left, anchor=north west, text width=28mm] at (135,-39) {local output port: VA needs a free RX slot of the class and reserves it};
  \node[stp, text width=25mm] (srv) at (49,-41) {\code{peekMessage}, \code{rx\_serving priority}: RESPONSE, then FORWARD, then REQUEST};
  \draw[arr] (87.3,-49) -- (srv.east |- 0,-49);
  \node[stp] (a3) at (2,-41) {\textbf{ADMIT-RF}: responses and forwards are always admitted to the processing queue. Requests wait for room.};
  \draw[arr] (srv.west |- 0,-49) -- (a3.east |- 0,-49);
  % ---------------- table ----------------
  \node[ttl] at (0,-74) {Slots one transition reserves (\code{CacheController::txNeed})};
  \node[box, anchor=north west, text width=160mm] at (0,-78) {\begin{tabular}{@{}p{26mm}p{48mm}p{70mm}@{}}
    \textbf{at} & \textbf{message popped} & \textbf{slots reserved}\\
    home (LLC bank) & demand request (GetS, GetM), hit or miss & 1 response and $\textit{caches}-1$ forwards (MOESI: \textit{caches}), plus 1 forward toward memory, reserved on every demand\\
    & PutS, PutM, PutE, PutO & 1 forward (the PutAck)\\
    & replacement & \textit{caches} forwards (one Inv each), plus 1 response toward memory\\
    & data from an L1 & 1 response toward memory (a possible write-back)\\
    cache (L1) & demand request (hits included) or replacement & 1 request\\
    & anything from the home without data & 2 responses (data to the requester and to the home)\\
  \end{tabular}\\[0.5ex]
  A refused message stays in the processing queue and is retried, like a full MSHR. The smallest workable \code{inject\_queue\_depth} is the number of caches (a replacement's fan-out). Buses report unlimited room, so admission never refuses on them.};
\end{tikzpicture}
\end{adjustbox}
\caption{The NI between a controller and its router. Transmit: admission reserves the injection
slots a transition will use, and the local port injects one flit per tick, oldest eligible packet
first. Receive:
a packet's RX slot is reserved at its last router, and the controller always admits responses and
forwards.}
\label{fig:noc-ni}
\end{figure}

\paragraph{Injection.} \code{pushMessage} classifies each destination's packet and puts it in the
injection queue of its class, \key{inject\_queue\_depth} packets deep (16 by default, 32 in the
grid presets). Each tick a local port sends at most one flit. Its candidates are the next flit of
every started packet that holds a credit and the head of every class queue that finds a free VC
of its class with a credit. The port's arbiter elects one. Under \code{FCFSArbiter} that is the
oldest eligible packet over all classes. A class head that finds no free VC with a credit is not a
candidate, so a younger packet of another class may leave first while an older one is blocked.
Within a class an endpoint never reorders its own messages. A started packet
keeps its VC until its tail has left, and no other packet's flits enter that VC in between
(wormhole). A new packet takes the lowest-index free VC of its class that its first route admits.
The local link costs nothing: a flit injected in a tick is in the router's input VC in the same
tick.

\paragraph{Ejection.} A packet's RX slot is reserved when its head wins VC allocation on the last
router's local output port. A local output port has no credit counter. This reservation plays
its role: a head wins a VC there only if the endpoint has a free RX slot in its class
(\code{canReserve}). When the tail leaves through the local port the packet is
complete and the message enters the RX queue of its class, \key{eject\_queue\_depth} messages
deep. The last router holds the message from the head's arrival to the tail's ejection, the one
store-and-forward step a delivered message object implies.

\paragraph{RX serving.} \key{rx\_serving} \code{priority} (the default) makes \code{peekMessage}
return a RESPONSE if there is one, else a FORWARD, else a REQUEST. \code{alternate} rotates over the
non-empty class queues and moves on every peek, as the legacy \code{MeshInterface} does.

\paragraph{Transmit-side admission (ADMIT-TX).} A push into an injection queue never fails. A
controller pops a message from its processing queue only if the injection queues its transition
will use have room (\code{txFreeSlots}), and reserves those slots before the actions run
(\cref{fig:noc-ni}, table). A refused message stays in the processing queue and is retried, as a
message is when the MSHR is full. The queues it waits for drain on their own, so no cycle closes.
A push refused despite the admission is an accounting error: the controller prints the queue and
the message and exits with status 2. The buses report unlimited room, so admission changes nothing
for them. The smallest workable \key{inject\_queue\_depth} is the number of caches, the fan-out of
the home's replacement. Below it the home can never admit a replacement and the run wedges. With
\env{OCTOPUS\_HANG\_DUMP} set, the watchdog then ends it with status 2.

\paragraph{Endpoint admission (ADMIT-RF).} The network is free of buffer deadlock only if every
endpoint eventually consumes any RESPONSE or FORWARD packet in its RX, whether or not it can take
a REQUEST. The cache controllers do: responses and forwards are always admitted to the processing
queue, and only requests wait for room. The memory controller \code{MainMemoryController} is the
exception. It accepts a \code{MEM\_READ}, a FORWARD, only with room in its queue
(\cref{sec:noc-banking}). It still drains, because it never waits on the network to serve its
queue: its fills go to an unbounded output queue.
\code{rx\_serving priority} lets the NI honour the property when all three RX queues hold packets. A unit test checks the property on the directory L1, the
directory LLC and the snooping L1 (\cref{sec:noc-validation}).

\section{The router}\label{sec:noc-router}

\begin{figure}[htbp]\centering
\begin{adjustbox}{max width=\textwidth}% Router microarchitecture and the per-tick phase order of Network::step (Network.cpp).
\begin{tikzpicture}[font=\sffamily\scriptsize, >=Latex, x=1mm, y=1mm,
  port/.style={draw=octrule, rounded corners=2pt, fill=octpaper, inner sep=0},
  q/.style={draw=octink, fill=white, minimum height=3.2mm, minimum width=3.2mm, inner sep=0},
  qh/.style={q, fill=octquiet!45},
  ctl/.style={draw=octpurple, fill=octio, thick, rounded corners=2pt, minimum height=11mm, text width=25mm, align=center, inner sep=2pt},
  cnt/.style={draw=octink, fill=octmem, minimum height=3.2mm, minimum width=5.5mm, inner sep=0, font=\sffamily\tiny},
  arr/.style={->, thick, octink},
  cred/.style={->, thick, octteal, densely dashed},
  ctlarr/.style={->, octpurple, semithick},
  lbl/.style={text=octquiet},
  ph/.style={draw=octink, fill=octpaper, rounded corners=2pt, align=left, inner sep=2.5pt, anchor=north west},
  ttl/.style={font=\sffamily\small\bfseries, text=octink, anchor=west}]
  % ---------------- control plane ----------------
  \node[ctl] (rc) at (30,30) {\textbf{RC} route computation\\\code{nextHop} gives \code{out\_port} and \code{vc\_mask}};
  \node[ctl] (va) at (75,30) {\textbf{VA} VC allocation\\one arbiter per output VC};
  \node[ctl] (sa) at (120,30) {\textbf{SA} switch allocation\\one arbiter per output port};
  \draw[ctlarr] (rc) -- node[above, lbl, align=center] {once per\\packet} (va);
  \draw[ctlarr] (va) -- node[above, lbl, align=center] {then once\\per flit} (sa);
  % ---------------- input ports ----------------
  \node[ttl] at (0,12) {input ports};
  \foreach \nm/\yy/\src in {link port 0/8/from upstream, local port 2/-22/from the NI} {
    \node[port, minimum width=42mm, minimum height=21mm, anchor=north west] at (0,\yy) {};
    \node[lbl, anchor=north west] at (1,\yy-0.5) {\textbf{\nm} (\src)};
    \foreach \c/\cn/\dy in {0/REQ/6, 1/FWD/11, 2/RSP/16} {
      \node[lbl, anchor=west, font=\sffamily\tiny] at (1,\yy-\dy) {\cn};
      \foreach \k in {0,1,2} \node[q] at (12+\k*3.4,\yy-\dy) {};
      \node[qh] at (22.2,\yy-\dy) {};
      \draw[arr] (23.8,\yy-\dy) -- (60,\yy-\dy);
    }
  }
  \node[lbl, align=left, anchor=north west] at (0,-45) {one FIFO per \code{[class][vc]}, \code{vc\_depth} flits each\\(shown: 3 classes, 1 VC). Grey: the flit at the front.};
  % ---------------- crossbar ----------------
  \draw[octink, thick, fill=octmem] (60,8) rectangle (90,-43);
  \node[font=\sffamily\scriptsize\bfseries] at (75,10.5) {crossbar};
  \foreach \yy in {2,-3,-8,-28,-33,-38} \draw[octquiet] (60,\yy) -- (90,\yy);
  \foreach \xx in {65,70,75,80,85} \draw[octquiet] (\xx,8) -- (\xx,-43);
  \node[lbl, align=center, fill=octmem, inner sep=1pt] at (75,-18) {one flit per input port\\and per output port\\per tick};
  % ---------------- output ports ----------------
  \node[ttl] at (122,11.5) {output ports};
  \node[port, minimum width=38mm, minimum height=21mm, anchor=north west] at (98,8) {};
  \node[lbl, anchor=north west] at (99,7.5) {\textbf{link port 0} (downstream)};
  \foreach \c/\cn/\dy in {0/REQ/6, 1/FWD/11, 2/RSP/16} {
    \node[lbl, anchor=west, font=\sffamily\tiny] at (99,8-\dy) {\cn};
    \node[cnt] at (113,8-\dy) {4};
    \node[lbl, anchor=west, font=\sffamily\tiny] at (117,8-\dy) {credits, holder};
    \draw[arr] (90,8-\dy) -- (98,8-\dy);
  }
  \node[port, minimum width=38mm, minimum height=21mm, anchor=north west] at (98,-22) {};
  \node[lbl, anchor=north west] at (99,-22.5) {\textbf{local port 2} (to the NI)};
  \foreach \c/\cn/\dy in {0/REQ/6, 1/FWD/11, 2/RSP/16} {
    \node[lbl, anchor=west, font=\sffamily\tiny] at (99,-22-\dy) {\cn};
    \node[cnt, fill=octctl] at (112,-22-\dy) {RX};
    \node[lbl, anchor=west, font=\sffamily\tiny] at (115.2,-22-\dy) {RX slot per packet};
    \draw[arr] (90,-22-\dy) -- (98,-22-\dy);
  }
  % links out
  \draw[arr] (136,1) -- (144,1) node[right, lbl, align=left] {link,\\\code{link\_latency}};
  \draw[arr] (136,-33) -- (144,-33) node[right, lbl, align=left] {RX queue\\of the class};
  % credits
  \draw[cred] (144,-9) -- (136,-9);
  \node[lbl, anchor=north west, text=octteal, align=left] at (137,-10) {credit back after\\\code{credit\_latency}};
  \draw[cred] (0,-3) -- (-8,-3);
  \node[lbl, text=octteal, align=right, anchor=east] at (-8.5,-3) {credit\\upstream\\when a flit\\leaves};
  \draw[cred] (0,-33) -- (-8,-33);
  \node[lbl, text=octteal, align=right, anchor=east] at (-8.5,-33) {credit\\to the NI};
  % control to data
  \fill[octpurple] (42,2) circle (0.5);
  \draw[ctlarr, densely dotted] (42,2) -- (42,14) -| (rc.south);
  \draw[ctlarr, densely dotted] (va.south) |- (101,17) -- (101,8);
  \draw[ctlarr, densely dotted] (sa.south) |- (84,13) -- (84,8);
  % ---------------- phases ----------------
  % every phase box has the same fixed height, text from the top
  \def\phbox#1{\parbox[c][22mm][t]{28mm}{\raggedright #1}}
  \node[ttl] at (-24,-56) {Each tick, over every router in ascending id and every port and VC in ascending index:};
  \node[ph] (p1) at (-24,-60) {\phbox{\textbf{1 Injection}\\each local port takes one flit from its NIs, with the credits seen at the start of the tick}};
  \node[ph] (p2) at (11,-60) {\phbox{\textbf{2 Arrivals}\\flits whose link time is due enter their VC. Due credits are added to the counters}};
  \node[ph] (p3) at (46,-60) {\phbox{\textbf{3 Route}\\every head at a FIFO front without a route gets \code{nextHop}}};
  \node[ph] (p4) at (81,-60) {\phbox{\textbf{4 VC allocation}\\heads that waited \code{router\_latency} compete. A local port also reserves the RX slot}};
  \node[ph] (p5) at (116,-60) {\phbox{\textbf{5 Switch and link}\\one flit per output port, given a credit and a free link. A tail sent on a local port is ejected}};
  \foreach \a/\b in {p1/p2, p2/p3, p3/p4, p4/p5} \draw[arr] (\a.east) -- (\b.west);
\end{tikzpicture}
\end{adjustbox}
\caption{A router. Each input port has one FIFO per class and VC. Route computation runs once per
packet, VC allocation claims an output VC, switch allocation moves one flit per output port per
tick. Credits flow back upstream when a flit leaves an input FIFO. Bottom: the five phases of
every tick.}
\label{fig:noc-router}
\end{figure}

A router has an input side and an output side per port (\cref{fig:noc-router}). An input port
holds \key{num\_classes} $\times$ \key{vcs\_per\_class} FIFOs of \key{vc\_depth} flits. An output
port keeps, per class and VC, a credit counter for the downstream FIFO and the packet that holds
the output VC. A VC may contain the tail of one packet followed by the head of the next, because
the upstream router releases its output VC when it sends the tail. The packet at the front is the
one being served.

A packet passes four stages:

\begin{description}
\item[RC, route computation] once per packet, when its head reaches the front of its FIFO. The
  routing algorithm returns an output port and a mask of the output VCs the packet may use.
\item[VA, VC allocation] once per packet. A head competes for a free output VC of its class on
  its output port no earlier than \key{router\_latency} ticks after it arrived (2 by default).
  Implicit routers of a \code{custom} graph add no latency. On a local output port the head also
  needs and reserves an RX slot. An algorithm with an escape network gets a second VA pass in the
  same tick, in which a head still without a VC may take its fallback route
  (\cref{sec:noc-routing}).
\item[SA, switch allocation] once per flit. A flit competes for its output port if it is at the
  front of its FIFO, its output VC has a credit and the link is free. A head also waits
  \key{switch\_latency} ticks after it won VA (0 by default, so VA and SA can fall in one tick).
  One flit leaves each input port and each output port per tick. Body and tail flits can leave
  in the tick they arrive.
\item[LT, link traversal] the flit is at the downstream router \key{link\_latency} ticks later
  (1 by default). A tail sent on a local port completes the packet.
\end{description}

\paragraph{Router options.} \key{router\_latency[<router id>](i)} overrides
\key{router\_latency} at one router, for a slower router in an otherwise uniform network. An id
that is not a router stops the run. \key{switch\_latency} adds a separate SA stage after VA at
every explicit router. \key{sa\_allocator} chooses the switch allocator. \code{output\_first}, the
default, lets every eligible input VC compete at its output port, and an input port that has sent
a flit drops out of the later elections of that tick. \code{input\_first} is a separable
input-first allocator. Each input port first picks one of its ready VCs with its own arbiter, and
the output ports then elect among at most one candidate per input port. Both allocators use
\key{arbiter\_type} (\cref{sec:noc-arbitration}). Implicit routers apply neither latency.

The network executes each tick in five phases, over all routers in ascending id and all ports and
VCs in ascending index: injection, arrivals (flits and credits whose time has come), route
computation, VC allocation, and switch allocation with link traversal and ejection. Injection uses
the credits visible at the start of the tick. Arrivals precede allocation, so a flit or a credit
that arrives in a tick is usable in that same tick. No decision depends on pointer values,
unordered containers or random numbers, so two runs of the same configuration produce
byte-identical reports.

\section{Flow control and latency}\label{sec:noc-flow}

\begin{figure}[!tbp]\centering
\begin{adjustbox}{max width=\textwidth}% Credit-based flow control between two routers, and the unloaded latency of a packet.
\begin{tikzpicture}[font=\sffamily\scriptsize, >=Latex, x=1mm, y=1mm,
  rtr/.style={draw=octpurple, fill=octio, very thick, rounded corners=2pt},
  q/.style={draw=octink, fill=white, minimum height=4.5mm, minimum width=4.5mm, inner sep=0},
  qf/.style={q, fill=octquiet!45},
  cnt/.style={draw=octink, fill=octmem, minimum height=6mm, minimum width=9mm, inner sep=0, font=\sffamily\small\bfseries},
  arr/.style={->, thick, octink},
  cred/.style={->, thick, octteal, densely dashed},
  lbl/.style={text=octquiet},
  box/.style={draw=octrule, rounded corners=2pt, fill=octpaper, inner sep=3pt, align=left},
  ttl/.style={font=\sffamily\small\bfseries, text=octink, anchor=west},
  hop/.style={draw=octpurple, fill=octio, thick, rounded corners=2pt, minimum width=27mm, minimum height=9mm, align=center, inner sep=1pt},
  nib/.style={draw=octink, fill=octquiet!25, minimum width=8mm, minimum height=6mm, inner sep=0, font=\sffamily\scriptsize\bfseries}]
  % ================= credit loop =================
  \node[ttl] at (0,25) {1. The credit loop of one VC};
  \draw[rtr] (0,18) rectangle (38,-6);   \node[font=\sffamily\small\bfseries, text=octpurple, anchor=north west] at (1,17) {router A};
  \draw[rtr] (86,18) rectangle (150,-6); \node[font=\sffamily\small\bfseries, text=octpurple, anchor=north west] at (87,17) {router B};
  % A: output VC credit counter
  \node[cnt] (cA) at (28,8) {3};
  \node[lbl, anchor=east, align=right] at (22,6) {credits of the\\output VC,\\$-1$ per flit sent};
  % link
  \draw[arr] (38,8) -- node[above, lbl] {flit, \code{link\_latency} $L$ ticks} (86,8);
  \node[qf, minimum width=6mm] at (62,4) {};
  % B: input VC
  \foreach \k in {0,1,2,3} \node[q] (b\k) at (100+\k*4.5,8) {};
  \node[qf] at (113.5,8) {};
  \node[lbl, anchor=north] at (107,4.5) {input VC FIFO, \code{vc\_depth} = 4};
  \draw[arr] (116,8) -- (146,8) node[midway, above, lbl, align=center] {forwarded (SA)\\slot freed};
  % credit return
  \draw[cred] (131,6) -- (131,-14) -- (28,-14) -- (cA.south);
  \node[lbl, text=octteal, fill=white, inner sep=1pt] at (78,-14) {credit, \code{credit\_latency} $c$ ticks after the flit leaves B's FIFO};
  % loop box
  \node[box, anchor=north west, text width=64mm] at (0,-20) {\textbf{Credit loop}: a slot is reusable $L + c$ ticks after its flit was sent.\\
    Per-VC throughput $\le \min\left(1, \dfrac{\textit{vc\_depth}}{L + c}\right)$ flits per tick.\\
    Measured with \code{router\_latency} 0 and $L = c = 1$: 0.5, 1 and 1 flit per tick for depths 1, 2 and 4.};
  \node[box, anchor=north west, text width=74mm] at (72,-20) {A flit is sent only when its downstream VC holds a credit. Credits return when B pops the flit, not when it arrives.
    \code{vc\_depth} smaller than a packet is wormhole switching. \code{vc\_depth} $\ge$ every packet's flit count is virtual cut-through.
    A local output port counts packets instead: one RX slot of the class, reserved at VC allocation.};
  % ================= unloaded latency =================
  \node[ttl] at (0,-58) {2. Unloaded latency of a packet of $f$ flits through $R$ routers};
  \node[nib] (s) at (4,-70) {NI};
  \node[hop] (h1) at (28,-70) {router 1\\\code{router\_latency}};
  \node[hop] (h2) at (68,-70) {router 2\\\code{router\_latency}};
  \node[lbl] (dots) at (92,-70) {$\cdots$};
  \node[hop] (h3) at (116,-70) {router $R$\\\code{router\_latency}};
  \node[nib] (d) at (146,-70) {NI};
  \draw[arr] (s) -- node[below, lbl] {0} (h1);
  \draw[arr] (h1) -- node[above, lbl] {link $L$} (h2);
  \draw[arr] (h2) -- (dots);
  \draw[arr] (dots) -- node[above, lbl] {link $L$} (h3);
  \draw[arr] (h3) -- node[below, lbl] {0} (d);
  \node[lbl, anchor=north west] at (0,-75.5) {local links cost nothing: an injected flit is in router 1 in the same tick};
  \node[box, anchor=north west, text width=146mm] at (0,-82) {\textbf{Unloaded latency}, from the head's injection to the tail's ejection:\\[0.8ex]
    \hspace*{30mm}\small$T = R\cdot\textit{router\_latency} + (R-1)\cdot\textit{link\_latency} + (f-1)$ ticks\scriptsize\\[0.8ex]
    The $(f-1)$ term is the tail following the head one tick per flit (cut-through: the head leaves a router before its tail arrives).
    The formula needs credits never to run out on the way: $\textit{vc\_depth} \ge \min(f, L + c)$.
    Example, the 2$\times$2 grid corner to corner ($R = 3$, defaults 2 and 1): $3\cdot 2 + 2\cdot 1 + 0 = 8$ ticks for a request and $8 + 4 = 12$ for a 5-flit data message.};
\end{tikzpicture}
\end{adjustbox}
\caption{Top: the credit loop of one VC and the throughput bound it sets. Bottom: the unloaded
latency of a packet.}
\label{fig:noc-credits}
\end{figure}

A flit is sent on a link only when its downstream VC holds a credit. The credit counter starts at
\key{vc\_depth} and drops by one per flit sent. When the downstream router forwards the flit and
frees its slot, a credit returns \key{credit\_latency} ticks later (by default the
\key{link\_latency}). A slot is therefore reusable $L + c$ ticks after its flit was sent, where $L$
is the link latency and $c$ the credit latency, and one VC carries at most
$\min(1, \textit{vc\_depth}/(L+c))$ flits per tick (\cref{fig:noc-credits}). Body flits reach this
bound. A head also waits \key{router\_latency} at each router. A \key{vc\_depth} smaller than a
packet is wormhole switching: a blocked packet then holds buffers in several routers. A
\key{vc\_depth} of at least every packet's flit count is virtual cut-through. A longer buffer
residency, as in a router that holds a slot through switch allocation and traversal, is modelled
with a larger \key{credit\_latency}.

\paragraph{Unloaded latency.} A packet of $f$ flits that crosses $R$ routers, and so $R-1$ links,
alone in the network, takes
\[
T = R\cdot(\textit{router\_latency} + \textit{switch\_latency}) + (R-1)\cdot\textit{link\_latency} + (f-1)
\]
ticks from the head's injection to the tail's ejection. With per-router overrides the first term
is the sum of $\textit{router\_latency}_r + \textit{switch\_latency}$ over the routers crossed,
and an implicit router adds nothing to it. Local links cost nothing. The $f-1$ term
is the tail following the head one tick per flit, since a head leaves a router before its tail
has arrived (cut-through). The formula holds while the credits never run out on the way, which
needs $\textit{vc\_depth} \ge \min(f, L + c)$ at every link latency. The default depth of 4 meets
it whenever $L + c \le 4$.

\paragraph{Clock.} The network is a \code{ClockedObj} with \key{m\_clk\_period} 50 (from
\file{Mesh.csv}). It steps on falling edges, so one tick is 100 time units, one core cycle as
shipped. All latency keys are in ticks. Halving \key{m\_clk\_period} gives two ticks per core
cycle.

\paragraph{Link options.} \key{class\_link\_latency(vi)} gives each class its own link latency,
one entry per class. \key{link\_model} \code{pipelined} (the default) lets a link accept one flit
per tick. \code{serial} holds the link for its whole latency per flit and delivers the flit
$\max(1, L-1)$ ticks after it was sent, which is the transfer model of the legacy mesh links.
\key{half\_duplex} 1 makes the two links of a connection one resource: one flit at a time in either
direction. The configured arbiter elects it over both routers' candidates.

\section{Routing}\label{sec:noc-routing}

\begin{figure}[htbp]\centering
\begin{adjustbox}{max width=\textwidth}% Four of the routing algorithms of Routing.cpp and the ordering each one declares
% (the full list is the table tab:noc-routing in chapter 9).
\begin{tikzpicture}[font=\sffamily\scriptsize, >=Latex, x=1mm, y=1mm,
  sr/.style={draw=octpurple, fill=octio, semithick, minimum size=4.2mm, inner sep=0, font=\sffamily\tiny},
  srs/.style={sr, fill=octink, text=white, font=\sffamily\tiny\bfseries},
  pth/.style={->, very thick, octink},
  alt/.style={->, very thick, octamber, densely dashed},
  esc/.style={->, very thick, octteal, densely dotted},
  ttl/.style={font=\sffamily\small\bfseries, text=octink, anchor=west},
  lbl/.style={text=octquiet, align=left, anchor=north west, text width=50mm},
  ord/.style={font=\sffamily\scriptsize\bfseries, anchor=west},
  meshgrid/.pic={
    \foreach \yy in {0,1,2,3} \foreach \xx in {0,1,2} \draw[octquiet] (\xx*10,-\yy*10) -- (\xx*10+10,-\yy*10);
    \foreach \yy in {0,1,2} \foreach \xx in {0,1,2,3} \draw[octquiet] (\xx*10,-\yy*10) -- (\xx*10,-\yy*10-10);
    \foreach \yy in {0,1,2,3} \foreach \xx in {0,1,2,3} {\pgfmathtruncatemacro{\rr}{\yy*4+\xx}\node[sr] at (\xx*10,-\yy*10) {\rr};}
  }]
  % ================= XY =================
  \begin{scope}[shift={(0,0)}]
    \node[ttl] at (-3,9) {\code{XY}};
    \pic {meshgrid};
    \draw[pth] (2.4,-30) -- (7.6,-30); \draw[pth] (12.4,-30) -- (17.6,-30); \draw[pth] (22.4,-30) -- (27.6,-30);
    \draw[pth] (30,-27.6) -- (30,-22.4); \draw[pth] (30,-17.6) -- (30,-12.4); \draw[pth] (30,-7.6) -- (30,-2.4);
    \node[srs] at (0,-30) {12}; \node[srs] at (30,0) {3};
    \node[lbl] at (-3,-39.5) {X to the destination column, then Y (then Z on a 3-D lattice). Deterministic, deadlock-free by construction. Any lattice.};
    \node[ord, text=octteal] at (-3,-36.5) {ordering: pair per class};
  \end{scope}
  % ================= WestFirst =================
  \begin{scope}[shift={(55,0)}]
    \node[ttl] at (-3,9) {\code{WestFirst}};
    \pic {meshgrid};
    % 3 -> 12: west first, all the way
    \draw[pth] (27.6,0) -- (22.4,0); \draw[pth] (17.6,0) -- (12.4,0); \draw[pth] (7.6,0) -- (2.4,0);
    \draw[pth] (0,-2.4) -- (0,-7.6); \draw[pth] (0,-12.4) -- (0,-17.6); \draw[pth] (0,-22.4) -- (0,-27.6);
    % 12 -> 3: adaptive between east and north
    \draw[alt] (2.4,-29) -- (7.6,-29); \draw[alt] (10,-27.6) -- (10,-22.4); \draw[alt] (12.4,-20) -- (17.6,-20);
    \draw[alt] (20,-17.6) -- (20,-12.4); \draw[alt] (22.4,-10) -- (27.6,-10); \draw[alt] (30,-7.6) -- (30,-2.4);
    \node[srs] at (0,-30) {12}; \node[srs] at (30,0) {3};
    \node[lbl] at (-3,-39.5) {Black: 3 to 12 turns west first. Amber: 12 to 3 picks, at each router, the minimal output (E or N) with the most credits. \code{NorthLast}, \code{NegativeFirst} and \code{OddEven} work alike with other turn rules. \code{grid2d} only.};
    \node[ord, text=octred] at (-3,-36.5) {ordering: none (adaptive)};
  \end{scope}
  % ================= AdaptiveMinimal =================
  \begin{scope}[shift={(110,0)}]
    \node[ttl] at (-3,9) {\code{AdaptiveMinimal}};
    \fill[octamber!15] (-3,-33) rectangle (33,3);
    \pic {meshgrid};
    \draw[alt] (2.4,-30) -- (7.6,-30); \draw[alt] (10,-27.6) -- (10,-22.4); \draw[alt] (10,-17.6) -- (10,-12.4);
    \draw[alt] (12.4,-10) -- (17.6,-10);
    \draw[esc] (22.4,-10) -- (27.6,-10); \draw[esc] (30,-7.6) -- (30,-2.4);
    \node[srs] at (0,-30) {12}; \node[srs] at (30,0) {3};
    \node[lbl] at (-3,-39.5) {Amber: any minimal output, most credits on the adaptive VCs. Teal: the adaptive VC was not granted, so the escape route, \code{XY} on \code{escape\_vc}, kept to the end. Needs \code{vcs\_per\_class} $\ge 2$, $\ge 3$ with wraps. Any lattice.};
    \node[ord, text=octred] at (-3,-36.5) {ordering: none (adaptive)};
  \end{scope}
  % ================= Table =================
  \begin{scope}[shift={(0,-76)}]
    \node[ttl] at (-3,9) {\code{Table}};
    \node[sr, minimum size=5mm] (a) at (0,0) {A}; \node[sr, minimum size=5mm] (b) at (15,0) {B};
    \node[sr, minimum size=5mm] (c) at (0,-14) {C}; \node[sr, minimum size=5mm] (d) at (15,-14) {D};
    \node[sr, minimum size=5mm] (e) at (30,-7) {E};
    \foreach \u/\v in {a/b, a/c, b/d, c/d, b/e, d/e} \draw[octquiet] (\u) -- (\v);
    \draw[pth] (c) -- (d); \draw[pth] (d) -- (e);
    \node[lbl] at (-3,-24) {Shortest paths by hops, or by \code{link\_weight}. Ties go to the lowest port index. \code{routing\_via} (\code{dest,neighbour}) and \code{routing\_table} (\code{dest,port}) override cells. Any topology, refused if its channel dependency graph has a cycle.};
    \node[ord, text=octteal] at (-3,-21) {ordering: pair per class};
  \end{scope}
  % ================= torus dateline =================
  \begin{scope}[shift={(55,-76)}]
    \node[ttl] at (-3,9) {\code{XY} with wraps};
    \foreach \xx in {0,1,2,3} \node[sr, minimum size=5mm] (t\xx) at (\xx*10,-4) {\xx};
    \foreach \xx/\nx in {0/1,1/2,2/3} \draw[octquiet] (t\xx) -- (t\nx);
    \draw[octquiet] (t3.north) to[out=60, in=120, looseness=0.5] (t0.north);
    \node[text=octquiet] at (15,5.5) {wrap link};
    \draw[pth] (t1) -- (t2);
    \node[text=octink, anchor=west] at (-3,-13.5) {1 to 2: VC 0};
    \draw[alt] (t3.south) to[out=-60, in=-120, looseness=0.45] (t0.south);
    \node[text=octamber, anchor=west] at (17,-13.5) {3 to 0 wraps: VC 1};
    \node[lbl] at (-3,-24) {Torus or ring: the shorter way round each wrapped dimension. A leg that crosses the wrap link uses VC 1, others VC 0 (the dateline rule), chosen from source and destination. Needs \code{vcs\_per\_class} $\ge 2$.};
    \node[ord, text=octteal] at (-3,-21) {ordering: pair per class};
  \end{scope}
  % ================= summary =================
  \begin{scope}[shift={(110,-76)}]
    \node[ttl] at (-3,9) {What a protocol can rely on};
    \node[lbl, text width=48mm, text=octink] at (-3,4) {\code{requires\_ordering(s),}\\\code{pair\_per\_class}, set by the directory presets, accepts only an algorithm that keeps two packets of one class from one source to one destination in order:\\[0.6ex]
      \code{Table} or \code{XY}, \code{vcs\_per\_class} 1\\
      \code{XY} on a torus or a ring (the dateline fixes the VC)\\[0.6ex]
      Anything else is refused at start-up, including the turn models, \code{O1TURN}, \code{Valiant} and \code{ROMM}. No algorithm orders packets of different classes.};
  \end{scope}
\end{tikzpicture}
\end{adjustbox}
\caption{Four of the routing algorithms on a 4$\times$4 grid and the dateline of \code{XY} on a
wrapped dimension, with the ordering each declares. Black: deterministic paths. Amber: adaptive
choices by credits. Teal: the escape route of \code{AdaptiveMinimal}. \Cref{tab:noc-routing}
lists all of them.}
\label{fig:noc-routing}
\end{figure}

A routing algorithm is a plug-in chosen by name with \key{routing}. Given the current router, the
input port and VC, the source and destination routers, the class and the packet, it returns the
output port and a mask of the output VCs the packet may use. It may read the current router's
credit counts and nothing else. Ten algorithms ship (\cref{tab:noc-routing}).
\code{auto}, the default, selects \code{XY} on \code{grid2d}, \code{torus2d}, \code{mesh3d} and
\code{torus3d}, \code{XY} on a \code{ring} when \key{vcs\_per\_class} is at least 2, and
\code{Table} everywhere else. A construction that does not fit an algorithm stops the run with
the reason, for example \code{O1TURN} on a 3-D mesh or \code{XY} on a torus with one VC.

\begin{table}[htbp]\centering\small
\caption{The routing algorithms, where they run and the ordering each declares
(\code{Routing.cpp}). A lattice is \code{grid2d}, \code{torus2d}, \code{ring}, \code{mesh3d} or
\code{torus3d}. \emph{Wraps} are the wrap links of a torus or a ring.}\label{tab:noc-routing}
\begin{tabularx}{\textwidth}{@{}llYl@{}}
\toprule algorithm & kind & topologies and VCs & ordering \\ \midrule
\code{Table} & deterministic & any, if the channel dependency graph is acyclic & pair per class \\
\code{XY} & deterministic & any lattice. With wraps, \key{vcs\_per\_class} $\ge 2$ & pair per class \\
\code{WestFirst} & turn model & \code{grid2d} & none \\
\code{NorthLast} & turn model & \code{grid2d} & none \\
\code{NegativeFirst} & turn model & \code{grid2d} & none \\
\code{OddEven} & turn model & \code{grid2d} & none \\
\code{O1TURN} & oblivious & a 2-D lattice without wraps (\code{grid2d}), \key{vcs\_per\_class} $\ge 2$ & none \\
\code{Valiant} & oblivious & a lattice without wraps (\code{grid2d}, \code{mesh3d}), \key{vcs\_per\_class} $\ge 2$ & none \\
\code{ROMM} & oblivious & as \code{Valiant} & none \\
\code{AdaptiveMinimal} & adaptive & any lattice, \key{escape\_vc} set, \key{vcs\_per\_class} $\ge 2$, with wraps $\ge 3$ & none \\
\bottomrule
\end{tabularx}
\end{table}

\paragraph{Deterministic.}
\begin{description}
\item[\code{Table}] any topology. One next hop per router and destination router. By default the
  shortest paths by hop count, found by breadth-first search from every destination, ties to the
  lowest port index. \key{link\_weight[<router>](vi)} gives \code{neighbour,weight} pairs (weights
  at least 1, default 1). With any weight set, \code{Table} computes the shortest weighted paths
  instead (Dijkstra) and prints how many router and destination pairs had several equal-cost
  ports, which it also breaks to the lowest port. \key{routing\_via[<router>](vi)} then sets single
  entries as \code{dest,neighbour} pairs: the route to router \code{dest} leaves through the link
  to router \code{neighbour}. It does not depend on port numbering, so it is the preferred way to
  state a route. \key{routing\_table[<router>](vi)} is applied last, as \code{dest,port} pairs. A
  neighbour that is not linked, or a port the router does not have, stops the run. The table must
  pass the channel dependency check (\cref{sec:noc-deadlock}).
\item[\code{XY}] any lattice. Dimension-order routing: along dimension 0 (X) to the destination's
  coordinate, then dimension 1 (Y), then dimension 2 (Z). On a wrapped dimension it takes the
  shorter way round, the positive direction on a tie, and applies the dateline rule. The leg in a
  dimension uses VC 1 if the shorter way from the source's coordinate to the destination's crosses
  the wrap link, and VC 0 otherwise. A torus or a ring with \code{XY} therefore needs
  \key{vcs\_per\_class} of at least 2, and only VCs 0 and 1 carry traffic.
\end{description}

\paragraph{Turn models.} \code{grid2d} only. Each forbids enough turns to break every cycle of
channel dependencies, so it is deadlock-free at any VC count, and every route it allows is
minimal. At each router the head takes, among the allowed outputs, the one whose VCs of its class
hold the most credits, ties to the lowest port.
\begin{description}
\item[\code{WestFirst}] all westward hops first. Then any minimal mix of east, north and south.
\item[\code{NorthLast}] northward hops last, with no turn after them. Before them any minimal mix
  of the east or west hops and the south hops.
\item[\code{NegativeFirst}] the negative directions, west and north, first and adaptively among
  them. Then the positive ones, east and south, adaptively.
\item[\code{OddEven}] Chiu's minimal odd-even rule: no turn from east to north or south in an even
  column, and no turn from north or south to west in an odd column. The column is the $x$
  coordinate.
\end{description}

\paragraph{Oblivious.} \code{O1TURN}, \code{Valiant} and \code{ROMM} choose once per packet,
when it is created, independently of the load. The choice is drawn from a hash of the packet's
network sequence number and \key{routing\_seed}, so a rerun takes the same paths and another seed
takes others. They run on lattices without wraps and split each class's VCs into two sets,
A $= [0, v/2)$ and B $= [v/2, v)$ for $v$ = \key{vcs\_per\_class}. Each set carries only
dimension-order routes, and a packet moves from A to B at most once, so the channel dependency
graph is acyclic.
\begin{description}
\item[\code{O1TURN}] each packet takes XY on set A or YX on set B, each with probability one half.
  2-D only.
\item[\code{Valiant}] dimension order to an intermediate router drawn uniformly among all
  routers, on set A, then dimension order to the destination on set B. The paths are not minimal.
\item[\code{ROMM}] as \code{Valiant}, with the intermediate drawn in each dimension between the
  source's and the destination's coordinates, so every path is minimal.
\end{description}

\paragraph{Adaptive.}
\begin{description}
\item[\code{AdaptiveMinimal}] any lattice, with \key{escape\_vc} set. Minimal adaptive routing with
  an escape network (Duato). The minimal outputs are one per dimension still to cover. The head
  asks for the one with the most credits on the adaptive VCs, which are every VC but the escape
  VCs, and carries the escape route as its fallback. If VA cannot grant the adaptive output, its
  second pass may grant the escape route instead, so a waiting head can always leave on the escape
  network. A head that sees no adaptive credit on any minimal output asks for the escape route
  directly. A packet that arrived over a link on an escape VC stays on the escape network to its
  destination. The escape route is dimension order on \key{escape\_vc}. On a wrapped dimension it
  uses two escape VCs with a dateline that depends only on the current router: a hop whose
  remaining leg still crosses the wrap link uses \key{escape\_vc}$+1$, any other hop
  \key{escape\_vc}. A grid or a 3-D mesh needs \key{vcs\_per\_class} of at least 2, one adaptive
  VC and the escape VC. A torus or a ring needs at least 3 and \key{escape\_vc}$+1$ below
  \key{vcs\_per\_class}. At start-up the network checks that the escape network is free of
  cycles (\cref{sec:noc-deadlock}).
\end{description}

\paragraph{Ordering.} Two packets of one class from one source to one destination arrive in the
order they were injected when the path and the VC are fixed by source, destination and class:
\code{Table} and \code{XY} with \key{vcs\_per\_class} 1, and \code{XY} on a torus or a ring, where
the dateline fixes the VC. These two algorithms declare \code{PAIR\_PER\_CLASS}. The turn models,
the oblivious algorithms and \code{AdaptiveMinimal} declare \code{NONE}. No algorithm orders
packets of different classes. The directory presets set the system-level key
\key{requires\_ordering(s),pair\_per\_class}, and the network refuses at start-up an algorithm that
declares \code{NONE}, and more than one VC per class unless the algorithm fixes the VC per pair
(\code{XY} with wraps). Multicast is not routed as such: it is one packet per destination.

\section{Arbitration}\label{sec:noc-arbitration}

\begin{figure}[htbp]\centering
\begin{adjustbox}{max width=\textwidth}% Arbitration on the network: where electIndex is called, the Candidate record, and the two
% policies (FCFSArbiter, RRArbiter) on one example.
\begin{tikzpicture}[font=\sffamily\scriptsize, >=Latex, x=1mm, y=1mm,
  site/.style={draw=octpurple, fill=octio, thick, rounded corners=2pt, text width=40mm, align=left, inner sep=2.5pt, anchor=north west},
  cell/.style={draw=octrule, minimum height=5mm, minimum width=13mm, inner sep=0},
  hcell/.style={cell, fill=octquiet!25, font=\sffamily\scriptsize\bfseries},
  own/.style={circle, draw=octink, fill=white, minimum size=6mm, inner sep=0, font=\sffamily\scriptsize\bfseries},
  owc/.style={own, fill=octquiet!35},
  arr/.style={->, thick, octink},
  lbl/.style={text=octquiet},
  ttl/.style={font=\sffamily\small\bfseries, text=octink, anchor=west},
  win/.style={draw=octteal, very thick, rounded corners=1pt}]
  % ---- call sites ----
  \node[ttl] at (0,8) {Where an arbiter elects};
  \node[site] (s1) at (0,2) {\textbf{Injection}, one arbiter per local port: started packets and queue heads of every class and every NI on the port};
  \node[site] (s2) at (0,-14) {\textbf{VC allocation}, one arbiter per output VC: routed heads that waited \code{router\_latency}};
  \node[site] (s3) at (0,-30) {\textbf{Switch allocation}, one arbiter per output port (per connection with \code{half\_duplex}): flits holding an output VC and a credit};
  % ---- candidates ----
  \node[ttl] at (52,8) {\code{Candidate \{index, owner, seq\}}};
  \foreach \c/\t [count=\i from 0] in {index/index, owner/owner, seq/seq} \node[hcell] at (58+\i*13, 0) {\t};
  \foreach \ix/\ow/\sq [count=\r from 1] in {0/2/41, 1/0/57, 2/2/38, 3/10/60} {
    \node[cell] at (58, -\r*5) {\ix}; \node[cell] at (71, -\r*5) {\ow}; \node[cell] at (84, -\r*5) {\sq};
  }
  \node[lbl, align=left, anchor=north west, text width=45mm] at (51.5,-24) {\code{owner} = \code{msg.owner}, an endpoint id (the requesting core for its requests, supplies and fills).\\\code{seq} = the packet's network-wide injection number: lower is older.};
  \foreach \st in {s1,s2,s3} \draw[arr] (\st.east) -- (51.2, -7.5);
  \draw[arr] (91,-7.5) -- (100,-7.5) node[right, lbl, align=left] {\code{electIndex(tick, candidates)}\\returns the winner's \code{index}};
  % ---- FCFS ----
  \node[ttl] at (100,8) {};
  \node[ttl] at (100,-21) {\code{FCFSArbiter}};
  \node[lbl, align=left, anchor=north west, text width=60mm] at (100,-25) {lowest \code{seq} wins, ties by lowest \code{index}:\\winner \textbf{index 2} (seq 38).};
  \draw[win] (51.5,-17.5) rectangle (90.5,-12.5);
  \node[text=octteal, font=\sffamily\tiny\bfseries, anchor=west] at (91,-15) {FCFS};
  % ---- RR ----
  \node[ttl] at (100,-37) {\code{RRArbiter}};
  \foreach \o/\a [count=\i from 0] in {0/90, 1/18, 2/-54, 3/-126, 10/162} {
    \ifnum\i=4 \node[owc] (o\i) at ($(113,-57)+(\a:11)$) {\o}; \else
    \ifnum\i=0 \node[owc] (o\i) at ($(113,-57)+(\a:11)$) {\o}; \else
    \ifnum\i=2 \node[owc] (o\i) at ($(113,-57)+(\a:11)$) {\o}; \else
    \node[own] (o\i) at ($(113,-57)+(\a:11)$) {\o}; \fi\fi\fi
  }
  \draw[->, octquiet, thick] (o2) to[bend left=30] node[pos=0.5, below, lbl] {skip} (o3);
  \draw[->, octteal, very thick] (o3) to[bend left=30] node[pos=0.5, left, text=octteal, font=\sffamily\tiny\bfseries] {wins} (o4);
  \node[lbl, anchor=north] at (113,-73) {owner ids in rotation (grey: has a candidate)};
  \node[lbl, align=left, anchor=north west, text width=40mm] at (128,-42) {pointer was on owner 2 (the last winner). It advances: owner 3 has no candidate and is skipped, owner 10 has one: winner \textbf{index 3} (seq 60). The pointer rests on 10. An owner's own candidates go oldest first.};
  \draw[win, densely dashed] (51.5,-22.5) rectangle (90.5,-17.6);
  \node[text=octteal, font=\sffamily\tiny\bfseries, anchor=west] at (91,-20) {RR};
\end{tikzpicture}
\end{adjustbox}
\caption{An election on the network: the allocator lists its candidates, and the arbiter returns
the winner's index. The same four candidates under the two shipped policies.}
\label{fig:noc-arbitration}
\end{figure}

Three allocators elect on the network: the injection of each local port, the VC allocation of each
output VC, and the switch allocation of each output port (of each connection under
\key{half\_duplex}). Each presents its candidates to an arbiter as a list of
\code{Candidate \{index, owner, seq\}} records and calls
\code{electIndex(tick, candidates)}, which returns the winner's \code{index}
(\cref{fig:noc-arbitration}). \code{owner} is the message's owner, an endpoint id: the
requesting core for its requests, supplies and fills. \code{seq} is the packet's network-wide
injection number, so a lower \code{seq} is an older packet. A single candidate wins without an
election. \key{arbiter\_type} selects the policy for all three:

\begin{description}
\item[\code{FCFSArbiter}] (default) the lowest \code{seq} wins, ties by the lowest index.
\item[\code{RRArbiter}] a pointer rotates over the endpoint ids in ascending order. The first
  owner after the pointer that has a candidate wins with its oldest candidate, and the pointer
  rests on it. An owner without a candidate does not use a turn. It shares its rotation code with
  the bus arbiter of the same name (\cref{sec:arbiters}).
\end{description}

The losers of a switch election are counted: each losing packet records the owner that beat it,
which feeds the \col{NoC Ahead} columns (\cref{sec:noc-observe}). Under \code{FCFSArbiter} the
older of two eligible packets goes first. Two packets that share a link therefore do not
interleave while both have a flit ready with a credit. When the older one is blocked, the younger
one's flits pass it. Under \code{RRArbiter} they interleave by owner. \code{TDMArbiter} and \code{PriorityArbiter} are refused on the network
as deferred.

\section{The worked example}\label{sec:noc-example}

\begin{figure}[htbp]\centering
\begin{adjustbox}{max width=\textwidth}% The worked example on the 2x2 grid of MultiCoreSystem_Grid (router_latency 2, link_latency 1,
% credit_latency 1, vc_depth 4, XY): a GetS from L1 0 to the LLC on router 3, then the 5-flit data
% response back. Cells mark the tick a flit leaves a router (switch allocation and link traversal).
\begin{tikzpicture}[font=\sffamily\scriptsize, >=Latex, x=1mm, y=1mm,
  cl/.style={draw=octrule!60, minimum width=9mm, minimum height=6mm, inner sep=0},
  wt/.style={cl, fill=octquiet!18, text=octquiet},
  hd/.style={cl, fill=octink, text=white, font=\sffamily\scriptsize\bfseries},
  bd/.style={cl, fill=octteal!35, font=\sffamily\scriptsize\bfseries},
  tl/.style={cl, fill=octteal!75, text=white, font=\sffamily\scriptsize\bfseries},
  rowl/.style={anchor=east, font=\sffamily\scriptsize\bfseries, align=right},
  hop/.style={->, octamber, thick},
  lbl/.style={text=octquiet},
  ttl/.style={font=\sffamily\small\bfseries, text=octink, anchor=west}]
  % ---------- axis ----------
  \foreach \t in {0,...,12} \node[lbl] at (32+\t*9, 6) {\t};
  \node[lbl, anchor=east] at (27,6) {tick};
  % ================= request =================
  \node[ttl] at (0,13) {GetS, 1 flit: L1 0 on R0 to the LLC on R3, route R0 $\to$ R1 $\to$ R3};
  \foreach \nm/\yy in {R0 (inject)/0, R1/-7, R3 (eject)/-14} \node[rowl] at (27,\yy) {\nm};
  \foreach \t in {0,...,12} \foreach \yy in {0,-7,-14} \node[cl] at (32+\t*9,\yy) {};
  \foreach \t/\yy in {0/0, 1/0, 3/-7, 4/-7, 6/-14, 7/-14} \node[wt] at (32+\t*9,\yy) {rl};
  \node[hd] at (32+2*9,0) {H/T};  \node[hd] at (32+5*9,-7) {H/T};  \node[hd] at (32+8*9,-14) {H/T};
  \draw[hop] (32+2*9+3,-3) -- (32+3*9-3,-4);
  \draw[hop] (32+5*9+3,-10) -- (32+6*9-3,-11);
  \node[anchor=west, text=octink, align=left] at (32+12*9+6,-7) {ejected at 8:\\$3\cdot 2 + 2\cdot 1 + 0$};
  % ================= response =================
  \node[ttl] at (0,-25) {Data, 5 flits: the LLC on R3 to L1 0 on R0, route R3 $\to$ R2 $\to$ R0 (X first)};
  \foreach \nm/\yy in {R3 (inject)/-38, R2/-45, R0 (eject)/-52} \node[rowl] at (27,\yy) {\nm};
  \node[rowl, font=\sffamily\scriptsize, text=octquiet] at (27,-31) {NI injects};
  \foreach \t in {0,...,12} \foreach \yy in {-31,-38,-45,-52} \node[cl] at (32+\t*9,\yy) {};
  \foreach \t/\f in {0/H, 1/1, 2/2, 3/3, 4/T} \node[cl, text=octink] at (32+\t*9,-31) {\f};
  \foreach \t/\yy in {0/-38, 1/-38, 3/-45, 4/-45, 6/-52, 7/-52} \node[wt] at (32+\t*9,\yy) {rl};
  \foreach \row/\ts in {-38/2, -45/5, -52/8} {
    \node[hd] at (32+\ts*9,\row) {H};
    \foreach \k in {1,2,3} \node[bd] at (32+\ts*9+\k*9,\row) {\k};
    \node[tl] at (32+\ts*9+4*9,\row) {T};
  }
  \draw[hop] (32+2*9+3,-41) -- (32+3*9-3,-42);
  \draw[hop] (32+5*9+3,-48) -- (32+6*9-3,-49);
  \node[anchor=west, text=octink, align=left] at (32+12*9+6,-45) {tail ejected at 12:\\$3\cdot 2 + 2\cdot 1 + 4$};
  % legend
  \begin{scope}[shift={(32,-62)}]
    \node[wt, minimum width=6mm] at (0,0) {rl}; \node[lbl, anchor=west] at (4,0) {head waits \code{router\_latency} before VC allocation};
    \node[hd, minimum width=6mm] at (76,0) {H}; \node[lbl, anchor=west] at (80,0) {head};
    \node[bd, minimum width=6mm] at (93,0) {1}; \node[lbl, anchor=west] at (97,0) {body};
    \node[tl, minimum width=6mm] at (110,0) {T}; \node[lbl, anchor=west] at (114,0) {tail};
    \node[lbl, anchor=north west, text width=130mm] at (-30,-5) {A filled cell is the tick a flit wins switch allocation and leaves the router. It arrives one tick later at the next router (amber), where a body flit can leave in the tick it arrives. On the last router the flit leaves through the local port: that is ejection. Credits never run out here, because a slot freed at tick $t$ is usable again at $t+2$ and \code{vc\_depth} is 4.};
  \end{scope}
\end{tikzpicture}
\end{adjustbox}
\caption{A GetS and its data on \code{MultiCoreSystem\_Grid} with the default timing, tick by
tick. Each filled cell is a flit leaving a router. The latencies match the unloaded formula.}
\label{fig:noc-example}
\end{figure}

On \code{MultiCoreSystem\_Grid}, L1 0 on router 0 sends a GetS to the LLC on router 3
(\cref{fig:noc-example}). The request is one flit. \code{XY} sends it east to router 1, then
south to router 3, then out of the LLC's local port. It waits \key{router\_latency} 2 at each of
the three routers and crosses two links: $3\cdot 2 + 2\cdot 1 + 0 = 8$ ticks. The 72-byte data
response is five flits and takes a different path, west to router 2 and then north to router 0,
because \code{XY} goes along X first: $3\cdot 2 + 2\cdot 1 + 4 = 12$ ticks. Request and response
use disjoint links. With the LLC on router 0 (\code{-p "interconnect.node\_map[10](i)=0"}) the same
pair costs $1\cdot 2 + 0 + 0 = 2$ and $1\cdot 2 + 0 + 4 = 6$ ticks. The validation measures exactly
these figures (\cref{sec:noc-validation}).

\section{Message classes and deadlock freedom}\label{sec:noc-deadlock}

\begin{figure}[htbp]\centering
\begin{adjustbox}{max width=\textwidth}% The three message classes (virtual networks) of the directory protocols, the memory messages,
% and the consumption order that makes them deadlock-free.
\begin{tikzpicture}[font=\sffamily\scriptsize, >=Latex, x=1mm, y=1mm,
  act/.style={draw=octrule, fill=octmem, rounded corners=2pt, minimum width=24mm, minimum height=10mm, align=center, font=\sffamily\scriptsize\bfseries},
  home/.style={act, fill=octctl, draw=octamber},
  mcn/.style={act, fill=white, draw=octteal, thick},
  rq/.style={->, thick, octink},
  fw/.style={->, thick, octpurple, dashed},
  rs/.style={->, very thick, octteal, double=white, double distance=0.5pt},
  m/.style={font=\sffamily\scriptsize, inner sep=1pt},
  lbl/.style={text=octquiet},
  ttl/.style={font=\sffamily\small\bfseries, text=octink, anchor=west},
  cbox/.style={draw=octink, minimum width=24mm, minimum height=8mm, align=center, font=\sffamily\scriptsize\bfseries},
  box/.style={draw=octrule, rounded corners=2pt, fill=octpaper, inner sep=3pt, align=left}]
  % ---------------- actors ----------------
  \node[ttl] at (0,13) {Who sends what, in which class};
  \node[act] (req) at (14,0) {L1 requester};
  \node[home] (home) at (90,0) {home\\(LLC bank)};
  \node[act] (own) at (14,-44) {L1 owner\\or sharer};
  \node[mcn] (mc) at (90,-44) {memory\\controller};
  % L1 requester <-> home
  \draw[rq] ($(req.east)+(0,3)$) -- node[m, above] {GetS, GetM, PutS, PutM, PutE, PutO} ($(home.west)+(0,3)$);
  \draw[fw] ($(home.west)+(0,0)$) -- node[m, pos=0.5, fill=white] {PutAck (no data)} ($(req.east)+(0,0)$);
  \draw[rs] ($(home.west)+(0,-3.5)$) -- node[m, below=0.8mm] {data from the home} ($(req.east)+(0,-3.5)$);
  % home <-> owner/sharer (diagonals, drawn left to right so the sloped labels sit outside the pair)
  \draw[fw] ($(home.south west)+(0,0)$) -- coordinate[pos=0.5] (fwm) ($(own.north east)+(-3,0)$);
  \draw[rs] ($(own.east)+(0,0)$) -- coordinate[pos=0.5] (rsm) ($(home.south)+(-6,0)$);
  \node[m, text=octpurple, rotate=32, anchor=south] at ($(fwm)+(-1.2,1.6)$) {FwdGetS, FwdGetM, Inv};
  \node[m, text=octteal, rotate=34, anchor=north] at ($(rsm)+(1.4,-2.0)$) {data to the home};
  % owner/sharer -> requester
  \draw[rs] (own.north) -- node[m, left=1.2mm, align=right, text=octteal] {data from\\the owner,\\InvAck from\\a sharer} (req.south);
  % home <-> memory controller
  \draw[fw] ($(home.south)+(0,0)$) -- node[m, left, pos=0.55] {\code{MEM\_READ}} ($(mc.north)+(0,0)$);
  \draw[rs] ($(mc.north)+(6,0)$) -- node[m, right, align=left, pos=0.45] {\code{FILL},\\\code{MEM\_WACK}} ($(home.south)+(6,0)$);
  \draw[rs] ($(home.east)+(0,0)$) -- ++(12,0) |- ($(mc.east)+(0,0)$);
  \node[m, text=octteal, anchor=west] at (116,-22) {\code{MEM\_WRITE}};
  % legend
  \node[ttl] at (138,13) {Classes};
  \draw[rq] (138,4) -- (148,4); \node[anchor=west] at (149,4) {\textbf{REQUEST}, class 0};
  \draw[fw] (138,-2) -- (148,-2); \node[anchor=west] at (149,-2) {\textbf{FORWARD}, class 1};
  \draw[rs] (138,-8) -- (148,-8); \node[anchor=west] at (149,-8) {\textbf{RESPONSE}, class 2};
  \node[lbl, anchor=north west, text width=36mm, align=left] at (137,-13) {A home's acknowledgement without data travels with its forwards, so it cannot overtake an earlier Inv or Fwd to the same cache.\\[0.6ex] Memory messages are on the network only with \code{memory\_on\_network(i),1}.};
  % ---------------- consumption order ----------------
  \node[ttl] at (0,-58) {Why three classes cannot deadlock};
  \node[cbox, fill=white] (c0) at (14,-76) {REQUEST};
  \node[cbox, fill=octquiet!20] (c1) at (54,-76) {FORWARD};
  \node[cbox, fill=octquiet!45] (c2) at (94,-76) {RESPONSE};
  \draw[->, very thick] (c0) -- (c1);
  \draw[->, very thick] (c1) -- (c2);
  \node[lbl, align=center, anchor=north] at (34,-81) {forwards,\\memory reads};
  \node[lbl, align=center, anchor=north] at (74,-81) {data, acks,\\fills};
  \draw[->, very thick] (c0.north) to[out=20, in=160] node[above, lbl] {data} (c2.north);
  \draw[->, very thick] (c2.east) to[out=20, in=-20, looseness=6] node[right, lbl, align=left] {data,\\acks} (c2.south east);
  \node[box, anchor=north west, text width=46mm] at (126,-58) {A cell that consumes a class emits only higher classes. RESPONSE, the top class, may also emit itself: every endpoint consumes responses unconditionally (\mbox{ADMIT-RF}). Each class has its own buffers at every hop and its own RX queue, so a full REQUEST network never holds up the packets that drain it. With two classes a forward waits behind requests, and the check refuses the protocol.};
\end{tikzpicture}
\end{adjustbox}
\caption{The classes of the directory protocols' messages and of the memory messages, and the
consumption order that makes three classes deadlock-free.}
\label{fig:noc-classes}
\end{figure}

A directory protocol has a chain of dependencies: a home answers a request with forwards or data,
and a cache answers a forward with data (\cref{fig:noc-classes}). If a request and the forward it
causes share buffers, a network full of requests can block the forwards that would drain it, and
the system deadlocks. Three classes cut the chain. A cell of an FSM that consumes a class emits
only higher classes, and the top class, RESPONSE, is consumed unconditionally (ADMIT-RF). Each
class has its own buffers at every hop, so a full REQUEST network never blocks a FORWARD or a
RESPONSE. The network checks both halves of the argument at start-up (\cref{fig:noc-deadlock}): the
routing when it is built, the protocol once every controller exists.

\begin{figure}[htbp]\centering
\begin{adjustbox}{max width=\textwidth}% The two start-up deadlock checks: the channel dependency graph of a routing
% algorithm (Routing.cpp) and the protocol-and-network check of the FSM tables (ProtocolCheck.cpp).
\begin{tikzpicture}[font=\sffamily\scriptsize, >=Latex, x=1mm, y=1mm,
  sr/.style={draw=octpurple, fill=octio, thick, minimum size=6mm, inner sep=0, font=\sffamily\scriptsize\bfseries},
  ch/.style={draw=octink, fill=white, rounded corners=1.5pt, minimum width=11mm, minimum height=5mm, inner sep=0.5pt, font=\sffamily\tiny},
  dep/.style={->, very thick, octred},
  lk/.style={->, semithick, octquiet},
  pth/.style={->, thick, octink, densely dashed},
  lbl/.style={text=octquiet},
  ttl/.style={font=\sffamily\small\bfseries, text=octink, anchor=west},
  box/.style={draw=octrule, rounded corners=2pt, fill=octpaper, inner sep=3pt, align=left},
  hd/.style={font=\sffamily\scriptsize\bfseries, anchor=west},
  ok/.style={text=octteal, font=\sffamily\scriptsize\bfseries},
  bad/.style={text=octred, font=\sffamily\scriptsize\bfseries}]
  % ================= routing level =================
  \node[ttl] at (0,14) {1. Routing level: the channel dependency graph};
  \node[sr] (r0) at (8,0) {0}; \node[sr] (r1) at (32,0) {1};
  \node[sr] (r2) at (32,-24) {2}; \node[sr] (r3) at (8,-24) {3};
  \draw[lk] (r0) -- (r1); \draw[lk] (r1) -- (r2); \draw[lk] (r2) -- (r3); \draw[lk] (r3) -- (r0);
  \node[lbl] at (20,7.5) {4-router ring, \code{Table}};
  \draw[pth] ($(r0.south east)+(1,-1)$) -- ($(r1.south west)+(-1,-1)$) -- ($(r2.north west)+(-1,1)$);
  \node[lbl, align=center] at (17.5,-12) {two-hop paths:\\0 to 2 via 1,\\1 to 3 via 2, ...};
  % channels
  \node[ch] (c01) at (60,0) {0$\to$1, vc 0};
  \node[ch] (c12) at (80,-12) {1$\to$2, vc 0};
  \node[ch] (c23) at (60,-24) {2$\to$3, vc 0};
  \node[ch] (c30) at (40,-12) {3$\to$0, vc 0};
  \draw[dep] (c01) to[bend left=15] (c12); \draw[dep] (c12) to[bend left=15] (c23);
  \draw[dep] (c23) to[bend left=15] (c30); \draw[dep] (c30) to[bend left=15] (c01);
  \node[lbl, align=center] at (60,-12) {nodes: (link, VC)\\edges: a path uses\\one, then the next};
  \node[box, anchor=north west, text width=88mm] at (0,-32) {Built when the network is constructed, from every source and destination router, with each hop's \code{vc\_mask}. A cycle means a set of packets can each hold one channel and wait for the next. For \code{Table} a cycle is fatal and the message names it. \code{XY}, \code{WestFirst} and \code{AdaptiveMinimal} (through its \code{XY} escape VC) are deadlock-free by construction on \code{grid2d}. The computed \code{Table} on a ring of four or more routers has this cycle at any VC count, so it is refused. A hand-written \code{routing\_table} that breaks the cycle passes.};
  % ================= protocol level =================
  \node[ttl] at (96,14) {2. Protocol and network};
  \node[box, anchor=north west, text width=64mm] at (96,9) {For every FSM cell with actions, of every controller on the network: the class the event consumes (\code{consumedBy}) and the classes its actions emit (\code{emittedBy}), each classified exactly as the network classifies messages.\\[0.6ex]
    \textbf{Rule}: emitted class $>$ consumed class, or both RESPONSE.\\
    Cells fed by the CPU, the memory bus or the controller itself are not constrained.};
  \node[hd] at (96,-27) {cell};
  \node[hd] at (121,-27) {consumes};
  \node[hd] at (139,-27) {emits};
  \draw[octrule] (96,-29.5) -- (162,-29.5);
  \node[anchor=west] at (96,-33) {LLC (S, GetM)};          \node[anchor=west] at (121,-33) {REQUEST}; \node[anchor=west] at (139,-33) {Inv, data}; \node[ok] at (160,-33) {ok};
  \node[anchor=west] at (96,-38) {L1 (S, Inv)};            \node[anchor=west] at (121,-38) {FORWARD}; \node[anchor=west] at (139,-38) {InvAck};    \node[ok] at (160,-38) {ok};
  \node[anchor=west] at (96,-43) {LLC, last InvAck};       \node[anchor=west] at (121,-43) {RESPONSE}; \node[anchor=west] at (139,-43) {data};     \node[ok] at (160,-43) {ok};
  \node[anchor=west] at (96,-48) {MC (read)};               \node[anchor=west] at (121,-48) {FORWARD}; \node[anchor=west] at (139,-48) {\code{FILL}};  \node[ok] at (160,-48) {ok};
  \draw[octrule] (96,-50.5) -- (162,-50.5);
  \node[anchor=west, text=octquiet] at (96,-54) {two classes: L1 (S, Inv)};  \node[anchor=west] at (139,-54) {class 0 $\to$ 0}; \node[bad] at (160,-54) {no};
  \node[box, anchor=north west, text width=64mm] at (96,-59) {Shipped directory MSI with three classes: 273 cells, 60 emissions, edges REQUEST$\to$FORWARD, REQUEST$\to$RESPONSE, FORWARD$\to$RESPONSE and RESPONSE$\to$RESPONSE. The order holds. With two classes 36 cells violate it and the run stops with status 2, unless \code{deadlock\_check(s),warn}.};
\end{tikzpicture}
\end{adjustbox}
\caption{The two start-up checks. Left: a cycle in the channel dependency graph of
\code{Table} routing on a ring. Right: the protocol-and-network check over the FSM cells.}
\label{fig:noc-deadlock}
\end{figure}

\paragraph{Routing level.} The channel dependency graph has one node per (link, VC) and an edge
from one channel to the next whenever some packet may hold the first and then ask for the second.
To build it the network starts at every source router, for every destination router, and follows
every output the algorithm may use there, its \emph{candidates}, with every VC each hop's mask
allows: all allowed ports of a turn model, both orders of \code{O1TURN}, every intermediate of
\code{Valiant}, the adaptive outputs and the escape fallback of \code{AdaptiveMinimal}. The walk
does not stop at the destination router unless the route there is the local port, so a
\code{Valiant} packet that passes its destination on the way to its intermediate is followed. A
cycle stops the run and the message names the channels on it.

At start-up each algorithm answers whether it is deadlock-free on the topology. \code{Table} runs
this check. \code{Table} never routes a packet on the escape VC, so an \key{escape\_vc} does not
lift a refusal. The computed \code{Table} on a ring of four or more routers always has a cycle
and is refused at any VC count. \code{routing\_via} or \code{routing\_table} entries that break the
cycle pass, and \code{XY} with two VCs is the ring's deadlock-free routing. \code{AdaptiveMinimal}
checks its escape network: a packet may enter it at any router and then follows the escape route
to its destination, so the graph is built from every entry router and destination. With one escape
VC on a torus this graph has a cycle through the wrap link, which is why a wrapped lattice needs
the two escape VCs. \code{XY}, the turn models and the oblivious algorithms are deadlock-free by
construction. On a network of up to 64 routers the start-up check still builds the graph over
all their candidates and reports a cycle as an internal error. Larger networks skip this
cross-check, because its cost grows with the square of the router count (about 6\,s for
\code{Valiant} on 16$\times$16) and the construction argument does not depend on the size. The
step T unit tests run the same check on several grid shapes and VC counts.

\paragraph{Protocol level.} Once every controller exists, the network reads each protocol's FSM
table. For every cell with actions it takes the class of the message the event consumes and the
classes of the messages the actions emit, classified exactly as the network classifies them. The
handlers name these through \code{consumedBy(event)} and \code{emittedBy(action)}. A cell may
emit only a class higher than the one it consumes, or RESPONSE from RESPONSE. Cells fed by the CPU,
the memory bus or the controller itself are not constrained. The verdict is printed once:

\begin{lstlisting}[style=plain]
Protocol-network check (3 classes): 273 cells, 60 network emissions; edges REQUEST->FORWARD REQUEST->RESPONSE FORWARD->RESPONSE RESPONSE->RESPONSE; consumption order holds
\end{lstlisting}

A violation names the cell and exits with status 2. \key{deadlock\_check(s),warn} prints it and
runs anyway, which the legacy two-class preset does. With two classes the directory tables give 36
violations. Each is a cell that answers a message in its own class: an Inv answered by an InvAck,
a PutS by a PutAck, or a GetS forwarded as a FwdGetS. A protocol without class
annotations is reported and skipped. With memory on the network the check adds the memory
controller: a read is answered by a fill (FORWARD to RESPONSE) and a write by an acknowledgement
(RESPONSE to RESPONSE).

\paragraph{When something wedges anyway.} With \env{OCTOPUS\_HANG\_DUMP}$=N$ the network watches
itself: $N$ ticks without a flit moving while packets are in flight end the run. It then prints
every occupied VC (packet, message id, class, owner, flits, output port, output VC and its credits)
and every NI's queues, finalises the raw trace and exits with status 2. A routing deadlock reads as
a ring of held VCs with zero credits. The controllers' watchdog runs the same dump.

\section{LLC banks and memory on the network}\label{sec:noc-banking}

\begin{figure}[htbp]\centering
\begin{adjustbox}{max width=\textwidth}% LLC banking and memory controllers as network endpoints: the two address maps and the
% endpoints of MultiCoreSystem_Grid2B2M, with the network legs of one LLC miss.
\begin{tikzpicture}[font=\sffamily\scriptsize, >=Latex, x=1mm, y=1mm,
  ln/.style={draw=octink, minimum width=5mm, minimum height=5mm, inner sep=0, font=\sffamily\tiny},
  bA/.style={ln, fill=octctl},
  bB/.style={ln, fill=octamber!45},
  mA/.style={ln, fill=white},
  mB/.style={ln, fill=octteal!30},
  rtr/.style={draw=octpurple, fill=octio, very thick, minimum size=9mm, inner sep=0, font=\sffamily\scriptsize\bfseries},
  ep/.style={draw=octrule, fill=octmem, rounded corners=1.5pt, minimum width=11mm, minimum height=4.2mm, inner sep=0.5pt, font=\sffamily\tiny\bfseries},
  bank/.style={ep, fill=octctl, draw=octamber},
  mport/.style={ep, fill=octctl, draw=octamber, densely dashed},
  mc/.style={ep, fill=white, draw=octteal, thick},
  leg/.style={->, very thick, octink},
  lbl/.style={text=octquiet},
  ttl/.style={font=\sffamily\small\bfseries, text=octink, anchor=west},
  box/.style={draw=octrule, rounded corners=2pt, fill=octpaper, inner sep=3pt, align=left}]
  % ================= address maps =================
  \node[ttl] at (0,14) {Two static address maps};
  \node[box, anchor=north west, text width=80mm] at (0,9) {With $\textit{line} = \textit{addr}/\textit{block}$:\\[0.6ex]
    \textbf{home of a line}, the same map at every requester\\
    \hspace*{3mm}\code{home\_ids[(line >{}> home\_interleave\_shift) mod N]}\\[0.6ex]
    \textbf{memory controller of a line}, a map on each LLC bank\\
    \hspace*{3mm}\code{mem\_ids[(line >{}> mem\_interleave\_shift) mod M]}};
  \node[lbl, anchor=west] at (0,-19) {lines};
  \foreach \i in {0,...,7} {
    \pgfmathtruncatemacro{\odd}{mod(\i,2)}
    \ifnum\odd=0 \node[bA] at (12+\i*5.2,-19) {\i}; \else \node[bB] at (12+\i*5.2,-19) {\i}; \fi
  }
  \node[lbl, anchor=west] at (54,-19) {home shift 0:};
  \node[lbl, anchor=west] at (54,-22.5) {banks 10, 11, 10, 11, ...};
  \node[lbl, anchor=west] at (0,-28) {pages};
  \foreach \i/\t in {0/0..63, 1/64..127, 2/128..191, 3/192..255} {
    \pgfmathtruncatemacro{\odd}{mod(\i,2)}
    \ifnum\odd=0 \node[mA, minimum width=10mm] at (14.5+\i*10.4,-28) {\t}; \else \node[mB, minimum width=10mm] at (14.5+\i*10.4,-28) {\t}; \fi
  }
  \node[lbl, anchor=west] at (54,-28) {memory shift 6: 4 KiB};
  \node[lbl, anchor=west] at (54,-31.5) {pages to 100, 101, ...};
  \node[lbl, anchor=north west, text width=80mm] at (0,-36) {Shipped in \code{MultiCoreSystem\_Grid2B2M}: two banks interleaved by line, two controllers by page, so every bank talks to both controllers. Lines never move and banks never exchange messages. With one bank this is the single LLC of the other presets.};
  % ================= endpoints of Grid2B2M =================
  \node[ttl] at (90,14) {Endpoints of \code{MultiCoreSystem\_Grid2B2M}};
  \node[rtr] (R0) at (112,0) {R0}; \node[rtr] (R1) at (148,0) {R1};
  \node[rtr] (R2) at (112,-34) {R2}; \node[rtr] (R3) at (148,-34) {R3};
  \foreach \a/\b in {R0/R1, R0/R2, R1/R3, R2/R3} \draw[octquiet, thick] (\a) -- (\b);
  \node[ep, anchor=east] (l0) at (105,5) {L1 0};   \node[bank, anchor=east] (b10) at (105,0) {bank 10}; \node[mport, anchor=east] (p110) at (105,-5) {port 110};
  \node[ep, anchor=west] (l1) at (155,2.5) {L1 1};   \node[mc, anchor=west] (m100) at (155,-2.5) {MC 100};
  \node[ep, anchor=east] (l2) at (105,-31.5) {L1 2}; \node[mc, anchor=east] (m101) at (105,-36.5) {MC 101};
  \node[ep, anchor=west] (l3) at (155,-29) {L1 3}; \node[bank, anchor=west] (b11) at (155,-34) {bank 11}; \node[mport, anchor=west] (p111) at (155,-39) {port 111};
  \foreach \e/\r in {l0/R0, b10/R0, p110/R0, l1/R1, m100/R1, l2/R2, m101/R2, l3/R3, b11/R3, p111/R3} \draw[octink, thin] (\e) -- (\r);
  % one miss: legs along the links (XY)
  \draw[leg] (145,-6) -- node[left, text=octink] {1 GetS} (145,-28);
  \draw[leg, octteal] (151,-28) -- node[right, text=octteal, align=left] {4 data} (151,-6);
  \draw[leg, octpurple, dashed] (141,-37) -- node[below, text=octpurple] {2 \code{MEM\_READ}} (119,-37);
  \draw[leg, octteal] (119,-31) -- node[above, text=octteal] {3 \code{FILL}} (141,-31);
  \node[box, anchor=north west, text width=74mm] at (90,-46) {Line 0x11c0 (line 71, page 1) is homed at bank 11 and kept by controller 101. A miss from L1 1 crosses the grid four times: GetS to bank 11, the read from port 111 to MC 101, the fill back, the data to L1 1. A bank is two endpoints: its id faces the L1s and its memory port (\code{mem\_port\_id}) faces memory.};
\end{tikzpicture}
\end{adjustbox}
\caption{The home map and the memory map, and the four network legs of one LLC miss on
\code{MultiCoreSystem\_Grid2B2M}.}
\label{fig:noc-banking}
\end{figure}

\paragraph{LLC banks.} With \key{llc\_banks(i)} $N > 1$ the LLC is $N$ complete LLC controllers,
configured as \code{llc\_controller[0]} to \code{llc\_controller[N-1]}. Each bank has its own
directory FSM, MSHRs, write-back buffer, data array, processing queue, endpoint id, NI and local
port. A
static home map shared by every requester assigns each line to one bank
(\cref{fig:noc-banking}):
\[
\textit{home} = \textit{home\_ids}\big[(\textit{line} \gg \textit{home\_interleave\_shift}) \bmod N\big],
\qquad \textit{line} = \textit{addr}/\textit{block}.
\]
\key{home\_ids(vi)} and \key{home\_interleave\_shift(i)} are cache keys (absent means the single
\key{m\_shared\_memory\_id}). Shift 0 interleaves consecutive lines, shift 6 interleaves 4\,KiB
pages. Lines never move between banks and banks never exchange messages, so the protocol does not
change. Capacity is per bank: a preset that banks the same LLC divides \code{cache\_size},
\code{num\_mshr}, \code{pwb\_size} and \code{processing\_queue\_size} by $N$, as
\code{MultiCoreSystem\_Grid2B} does. Without \code{memory\_on\_network} the banks share one memory
bus, listed in \code{bus.m\_lower\_level\_ids} and arbitrated with its configured policy, and the
memory answers the bank that asked.

\paragraph{Memory controllers as endpoints.} \key{memory\_on\_network(i),1} makes every memory
controller a network endpoint with its own id, NI and local port. A bank then reaches memory through a
second endpoint, its memory port \key{mem\_port\_id}. The bank keeps its two interfaces, lower
toward the L1s and upper toward memory, and the system gives it two NIs. The memory controller
answers \code{from}, which is the memory port. Memory messages travel in the classes of
\cref{tab:noc-classes}: a read in FORWARD, writes, fills and acknowledgements in RESPONSE. That is
why the protocol check passes with the memory path included. \code{MainMemoryController} always
accepts a write, whose number the banks' write-back buffers bound. It accepts a read only with room
in its queue, and its responses wait in an unbounded ordered output queue instead of being refused.
\code{MCsimInterface} queues every message it receives.

\paragraph{Several controllers.} \key{mem\_ids(vi)} and \key{mem\_interleave\_shift(i)} on the
banks are a second map, from a line to a memory controller, of the same form as the home map. The
system builds one memory controller per id in the union of the banks' \key{mem\_ids}. Several
controllers need \key{memory\_on\_network}. The same bits as the home map give each bank its own
controllers. Different bits, as in \code{Grid2B2M}, let every bank use every controller.
\key{main\_memory\_type(s),MCsim} builds one MCsim instance per controller. Each one sees the line
address with the controller-select bits removed, so it decodes a dense address space. The number of
controllers must then be a power of two.

\begin{figure}[htbp]\centering
\begin{adjustbox}{max width=\textwidth}% The read-after-write rule to memory: an LLC bank keeps every line it writes until the memory
% controller acknowledges the write, and fills a read of that line from its own copy meanwhile.
\begin{tikzpicture}[font=\sffamily\scriptsize, >=Latex, x=1mm, y=1mm,
  hdr/.style={draw=octamber, fill=octctl, rounded corners=2pt, minimum width=30mm, minimum height=8mm, align=center, font=\sffamily\scriptsize\bfseries},
  mch/.style={hdr, draw=octteal, fill=white, thick},
  life/.style={octquiet, thick, densely dashed},
  wr/.style={->, very thick, octteal},
  ack/.style={->, thick, octteal, densely dotted},
  rd/.style={->, thick, octpurple, dashed},
  note/.style={draw=octrule, fill=octpaper, rounded corners=1.5pt, align=left, inner sep=2pt, text width=46mm},
  held/.style={fill=octamber!35, draw=octamber},
  lbl/.style={text=octquiet},
  m/.style={fill=white, inner sep=1pt}]
  \node[hdr] (b) at (57,0) {LLC bank\\(pending-write table)};
  \node[mch] (mc) at (127,0) {memory controller};
  \draw[life] (57,-4) -- (57,-78); \draw[life] (127,-4) -- (127,-78);
  % time arrow
  \draw[->, octquiet] (0,-6) -- (0,-76) node[midway, left, rotate=90, anchor=south] {time};
  % write
  \draw[held] (55.5,-10) rectangle (58.5,-44);
  \draw[wr] (57,-10) -- node[m, sloped, above] {\code{MEM\_WRITE} X, data (RESPONSE)} (127,-20);
  \node[note, anchor=east] at (54,-12) {evicts dirty line X: the write leaves, X and its data stay in the table (counted in \code{pwb\_size})};
  % read while held
  \node[note, anchor=east] at (54,-31) {a miss on X now is filled from the table (\code{FILL\_ROLLBACK}): no read is sent};
  % ack
  \draw[ack] (127,-22) -- node[m, sloped, above] {\code{MEM\_WACK}, header only, id 0 (RESPONSE)} (57,-44);
  \node[note, anchor=west, text width=36mm] at (130,-22) {acknowledges on \emph{acceptance}: from here the controller's queue orders X};
  \node[note, anchor=east] at (54,-46) {acknowledgement: X leaves the table};
  % later read
  \draw[rd] (57,-56) -- node[m, sloped, above] {\code{MEM\_READ} X (FORWARD)} (127,-62);
  \draw[wr] (127,-66) -- node[m, sloped, below] {\code{FILL} X (RESPONSE)} (57,-74);
  \node[note, anchor=west, text width=36mm] at (130,-63) {\code{MainMemoryController} serves in arrival order, so the write comes first. \code{MCsimInterface} answers a read of a line whose write MCsim has not performed from that write.};
  \node[lbl, anchor=north west, text width=160mm] at (0,-82) {The guarantee does not rest on network order: reads and writes travel in different classes and may reach different controllers. Reads and writes of different lines may reorder. \code{OCTOPUS\_MEM\_HAZARD=1} counts any read that overtakes an older write of its line (\code{=fatal} stops at the first).};
\end{tikzpicture}
\end{adjustbox}
\caption{Read after write to memory: the bank keeps a written line until the controller
acknowledges it, and a controller fills a read from a write it has not yet performed.}
\label{fig:noc-raw}
\end{figure}

\paragraph{Read after write.} A read of a line observes every earlier write of that line. Reads
and writes of different lines may reorder. The guarantee does not rest on network order, since a
read and a write travel in different classes and may reach different controllers
(\cref{fig:noc-raw}). A memory controller acknowledges a write when it accepts it, with a
header-only \code{MEM\_WACK}. Until then the bank keeps the line and its data, counted in its
\code{pwb\_size}, and fills a read of that line from its own copy without sending one. After the
acknowledgement the controller's queue orders the line: \code{MainMemoryController} serves in
arrival order, and \code{MCsimInterface} answers a read of a line whose write MCsim has not
performed from that write. \env{OCTOPUS\_MEM\_HAZARD}$=1$ counts any read that overtakes an older
write of its line, and \code{=fatal} stops at the first one.

\paragraph{Accounting.} Memory packets are logged as \code{MEM\_BUS}, ENTER at injection and EXIT
at ejection. The report's L2-DRAM bus stages then contain the network time between the memory port
and the controller, and the event-path tiling still sums to the total
(\cref{ch:request-end-to-end}).

\section{The legacy single-switch NoC as a configuration}\label{sec:noc-legacy}

\begin{figure}[htbp]\centering
\begin{adjustbox}{max width=\textwidth}% The legacy single-switch NoC expressed as a network (MultiCoreSystem_NetNoC.csv).
\begin{tikzpicture}[font=\sffamily\scriptsize, >=Latex, x=1mm, y=1mm,
  imp/.style={draw=octrule, fill=octmem, rounded corners=1.5pt, minimum width=11mm, minimum height=6mm, inner sep=0.5pt, align=center, font=\sffamily\scriptsize\bfseries},
  llc/.style={imp, fill=octctl, draw=octamber},
  sw/.style={draw=octpurple, fill=octio, very thick, minimum size=10mm, inner sep=0, align=center, font=\sffamily\scriptsize\bfseries},
  hd/.style={<->, semithick, octink},
  hdl/.style={{Latex[length=2.2mm,width=1.6mm]}-{Latex[length=2.2mm,width=1.6mm]}, semithick, octink, densely dashed, shorten <=0.4mm, shorten >=0.4mm},
  lbl/.style={text=octquiet},
  ttl/.style={font=\sffamily\small\bfseries, text=octink, anchor=west},
  box/.style={draw=octrule, rounded corners=2pt, fill=octpaper, inner sep=3pt, align=left}]
  \node[ttl] at (0,14) {\code{MultiCoreSystem\_NetNoC}: the legacy NoC on the network engine};
  \node[sw] (s) at (36,-38) {200\\[-0.4ex]\mdseries\tiny router};
  \foreach \i/\x in {0/6,1/26,2/46,3/66} \node[imp] (c\i) at (\x,-18) {L1 \i};
  \node[llc] (l) at (36,2) {LLC 10};
  \foreach \i in {0,1,2,3} \draw[hd] (c\i) -- (s);
  \foreach \i in {0,1,2,3} \draw[hdl] (c\i) -- (l);
  \node[lbl, anchor=north west, text width=80mm] at (0,-48) {Each box is an endpoint on an implicit router of its own (no \code{router\_latency}). Solid: through the explicit router 200 (\code{router\_latency 1}). Dashed: direct links to the LLC. Every connection is one half-duplex link, shared by requests (2 ticks) and data (5 ticks) and held for the whole latency (\code{link\_model serial}). The legacy \code{NoC} had a separate request link and data link per connection.};
  % recipe
  \node[box, anchor=north west, text width=70mm] at (86,9) {\textbf{The keys that reproduce it}\\[0.4ex]
    \code{topology custom}, the legacy \code{connection\_matrix}, \code{router\_ids 200}\\
    \code{routing Table}\\
    \code{flit\_bytes 0} (one flit per message)\\
    \code{router\_latency 1}\\
    \code{num\_classes 2} (request, data)\\
    \code{class\_link\_latency 2,5}\\
    \code{link\_model serial}, \code{half\_duplex 1}\\
    \code{vc\_depth 32}, \code{rx\_serving alternate}\\
    \code{inject\_queue\_depth 1024}\\
    \code{deadlock\_check warn}};
  % path costs
  \node[ttl] at (86,-37) {One message, ticks};
  \node[box, anchor=north west, text width=70mm] at (86,-41) {\begin{tabular}{@{}lcc@{}}
    path & \code{NoC} & \code{Network} \\
    L1 to LLC, request & 1 & 1\\
    L1 to LLC, data & 4 & 4\\
    L1 to L1 through 200, request & 3 & 3\\
    L1 to L1 through 200, data & 9 & 9\\
  \end{tabular}};
\end{tikzpicture}
\end{adjustbox}
\caption{\code{MultiCoreSystem\_NetNoC}: the legacy \code{NoC} of \cref{ch:interconnects} as a
custom network. Single messages take exactly the legacy number of ticks.}
\label{fig:noc-legacy}
\end{figure}

The legacy \code{NoC} is the degenerate case of the network: point-to-point links, each core
linked to the LLC and to one switch, half duplex, a message held on a link for its whole latency.
The legacy controller gives every connection two links, a request link and a data link. The network
has one link per direction, shared by both classes, and under half duplex one per connection. \code{MultiCoreSystem\_NetNoC} sets the network up that way
(\cref{fig:noc-legacy}). The legacy matrix becomes a \code{custom} graph with switch 200 as the
only explicit router, \code{Table} routing with \key{routing\_via} entries that send L1-to-L1
messages through the switch, one flit per message, two classes, serial half-duplex
links with the class latencies 2 and 5, \code{alternate} RX serving and the protocol check in
\code{warn} mode.

Every single-message path takes exactly its legacy number of ticks, and all ten EEMBC traces
complete with average latencies within 0.2 cycles and finish cycles within $-0.5\,\%$ to
$+0.1\,\%$ of the legacy \code{NoC}. The reports are not byte-identical, and worst cases differ:
the request-bus worst case is on average 10.7 cycles higher. Three modelled differences cause
it. An NI injects one flit per tick, where the legacy node sent on all its links in the same cycle.
One input VC per class makes a message whose link is busy block the messages behind it, where the
legacy node kept one queue per link. A data flit holds the shared link for 5 ticks, so a request
on the same connection waits behind it, where the legacy request link was free. All three are
modelled on purpose. The legacy
\code{Mesh} and \code{NoC} classes remain selectable with \code{interconnect\_type Mesh} and
\code{NoC}.

\section{Network-only mode}\label{sec:noc-generators}

\code{-s NetworkOnly} builds a network and one \code{TrafficGenerator} per endpoint, and no
caches. A generator is an endpoint that injects synthetic packets and consumes what it receives at
once, so it satisfies ADMIT-RF trivially. It is configured under \code{generator[<id>]} or
\code{generator[*]}, with defaults in \file{configuration/TrafficGenerator.csv}
(\cref{tab:noc-gen}). The run ends when every generator has used its budget and no packet is in
flight. The generators are closed-loop: a refused packet is retried, not dropped, so above
saturation the offered load limits itself. Saturation is therefore read from the accepted rate:
the first offered rate whose accepted rate falls below 0.9 of it.

\paragraph{Patterns.} Under \code{uniform} every packet goes to a random other endpoint. Under
\code{hotspot} every packet goes to \key{hotspot\_id}, and under \code{neighbour} to the next
endpoint id, cyclically. The permutation patterns give every generator one fixed partner, which
the system computes from the topology at start-up (\cref{tab:noc-patterns}). In the bit patterns
$i$ is the endpoint's index among the endpoint ids in ascending order, $N$ is the number of
endpoints, and $N$ must be a power of two. A permutation other than \code{transpose} whose partner
is the endpoint itself, or no endpoint, makes that generator send nothing. \code{transpose} keeps
its fallback to the first other endpoint.

\begin{table}[htbp]\centering\small
\caption{Permutation patterns (\code{NetworkOnly.cpp}).}\label{tab:noc-patterns}
\begin{tabularx}{\textwidth}{@{}lY@{}}
\toprule pattern & partner of an endpoint \\ \midrule
\code{transpose} & the first endpoint at the router with swapped coordinates, $(x,y) \to (y,x)$, on a 2-D lattice. Without one (a 3-D or non-lattice topology, a swapped router outside the grid or without an endpoint) the first other endpoint \\
\code{bitcomp} & index $i$ with all $\log_2 N$ bits complemented \\
\code{bitrev} & index $i$ with its $\log_2 N$ bits in reverse order \\
\code{shuffle} & index $i$ rotated left by one bit \\
\code{tornado} & on a lattice, the first endpoint at the router whose coordinate in every dimension is $c' = (c + \lceil k/2 \rceil - 1) \bmod k$. Otherwise index $(i + \lceil N/2 \rceil - 1) \bmod N$ \\
\code{randperm} & index $\pi(i)$ of one random cyclic permutation $\pi$ of $0..N-1$: Sattolo's algorithm over \code{std::mt19937} seeded with \key{perm\_seed}, the same in every generator that has the same \key{perm\_seed}. A single cycle has no fixed point, so every endpoint sends and receives \\
\bottomrule
\end{tabularx}
\end{table}

\paragraph{Injection.} \key{injection} \code{bernoulli}, the default, injects a packet in each
tick with probability \key{rate}/1000. \code{onoff} makes the generator a two-state Markov source
that starts off. In each tick it first switches state, off with probability
\key{on\_to\_off}/1000 when on and on with probability \key{off\_to\_on}/1000 when off, and then,
if on, injects with probability \key{rate}/1000. The long-run share of time on is
$\textit{off\_to\_on} / (\textit{on\_to\_off} + \textit{off\_to\_on})$, so the mean offered rate is
that share times \key{rate}, in bursts.

\begin{table}[htbp]\centering\small
\caption{Traffic generator keys.}\label{tab:noc-gen}
\begin{tabularx}{\textwidth}{@{}llY@{}}
\toprule key & default & meaning \\ \midrule
\key{pattern} & uniform & \code{uniform}, \code{hotspot}, \code{neighbour} (the next endpoint id), \code{none}, or a permutation: \code{transpose}, \code{bitcomp}, \code{bitrev}, \code{shuffle}, \code{tornado}, \code{randperm} (\cref{tab:noc-patterns}) \\
\key{rate} & 50 & packets per 1000 ticks, from 0 to 1000. Under \code{onoff} the rate while on \\
\key{injection} & bernoulli & \code{bernoulli} or \code{onoff} \\
\key{on\_to\_off} & 100 & \code{onoff}: per-1000 probability per tick of switching off, 0 to 1000 \\
\key{off\_to\_on} & 100 & \code{onoff}: per-1000 probability per tick of switching on, 0 to 1000 \\
\key{perm\_seed} & 1 & \code{randperm}: the seed of the permutation \\
\key{payload\_bytes} & 64 & payload per packet, so the flit count follows \cref{sec:noc-packets}. 0 is a header-only packet \\
\key{cls} & 2 & class to inject on: 0, 1 or 2 \\
\key{hotspot\_id} & 0 & destination of the \code{hotspot} pattern \\
\key{packets} & 1000 & packets to inject, $-1$ for no limit \\
\key{cycles} & $-1$ & last injecting cycle, $-1$ for none \\
\key{seed} & 1 & the generator's own random seed, combined with its id \\
\key{start} & 0 & first injecting cycle \\
\key{consume\_after} & 0 & cycles before the generator drains its RX, $-1$ for never (a stalled sink) \\
\key{multicast} & 1 & destinations per message \\
\key{untagged} & 0 & 1 leaves the message kind unknown (tests the classifier's fallback) \\
\bottomrule
\end{tabularx}
\end{table}

\section{What the network reports}\label{sec:noc-observe}

\begin{figure}[htbp]\centering
\begin{adjustbox}{max width=\textwidth}% What the network records for one packet: raw-trace records per router (Logger role NOC), the
% bus records at the NI, and the trackers written into the LatencyReport at ejection.
% Illustration: the GetS of the worked example, delayed 4 ticks at R1 by an older write-back.
\begin{tikzpicture}[font=\sffamily\scriptsize, >=Latex, x=1mm, y=1mm,
  lane/.style={draw=octrule!50, fill=octpaper},
  stay/.style={draw=octpurple, fill=octio, thick},
  lost/.style={draw=octred, fill=octred!20, thick},
  busi/.style={draw=octink, fill=octquiet!25, thick},
  mk/.style={circle, fill=octink, inner sep=0, minimum size=1.6mm},
  rowl/.style={anchor=east, font=\sffamily\scriptsize\bfseries, align=right},
  lbl/.style={text=octquiet},
  ttl/.style={font=\sffamily\small\bfseries, text=octink, anchor=west},
  box/.style={draw=octrule, rounded corners=2pt, fill=octpaper, inner sep=3pt, align=left}]
  \def\tk{8}
  % axis
  \foreach \t in {0,...,12} { \node[lbl] at (30+\t*\tk, 6) {\t}; \draw[octrule!25] (30+\t*\tk,3) -- (30+\t*\tk,-36); }
  \node[lbl, anchor=east] at (25,6) {tick};
  % lanes
  \foreach \nm/\yy in {{\code{REQ\_BUS}}/0, {\code{NOC} R0}/-9, {\code{NOC} R1}/-18, {\code{NOC} R3}/-27} {
    \draw[lane] (28,\yy+3) rectangle (30+12*\tk+4,\yy-3); \node[rowl] at (25,\yy) {\nm};
  }
  % bus record: inject (ENTER, trace only) to eject (EXIT, trace and report)
  \draw[busi] (30,2) rectangle (30+12*\tk,-2);
  \node at (30+6*\tk,0) {packet: ENTER at injection, EXIT at ejection};
  % router stays: ENTER .. EXIT, SERVICE marked
  \draw[stay] (30,-7) rectangle (30+2*\tk,-11);                  \node[mk] at (30+2*\tk,-9) {};
  \draw[stay] (30+3*\tk,-16) rectangle (30+5*\tk,-20);           \node[mk] at (30+5*\tk,-18) {};
  \draw[lost] (30+5*\tk,-16) rectangle (30+9*\tk,-20);
  \node[text=octred] at (30+7*\tk,-18) {loses SA 4 times};
  \draw[stay] (30+10*\tk,-25) rectangle (30+12*\tk,-29);         \node[mk] at (30+12*\tk,-27) {};
  \node[lbl, anchor=west] at (30+2*\tk+1,-9) {EXIT};
  \node[lbl, anchor=west] at (30+9*\tk+1,-18) {EXIT};
  \node[lbl, anchor=east] at (30+10*\tk-1,-27) {ENTER};
  % legend
  \node[mk] at (31,-40) {}; \node[lbl, anchor=west] at (33,-40) {SERVICE: the head wins VC allocation};
  \draw[lost] (88,-41.5) rectangle (94,-38.5); \node[lbl, anchor=west] at (95,-40) {ready but beaten at switch allocation};
  % trackers
  \node[ttl] at (0,-50) {Written at ejection into the request's \code{LatencyReport} row};
  \node[box, anchor=north west, text width=160mm] at (0,-54) {\begin{tabular}{@{}p{30mm}p{124mm}@{}}
    \col{NoC Hops} = 3 & routers entered, summed over the request's packets\\
    \col{NoC Transfer} = 8 & what the packets would take unloaded: $\textit{router\_latency}$ per router, link ticks, and $f-1$\\
    \col{NoC Wait} = 4 & the rest of injection-to-ejection: waiting for a VC, the switch or a credit. Wait + Transfer = 12\\
    \col{NoC Ahead} = 4 & flits of other packets that won a switch election this packet's flit lost\\
    \col{NoC Ahead Winner} = 1 & the owner that won most often ($-1$ if never beaten)\\
  \end{tabular}\\[0.4ex]
  With \code{memory\_on\_network} the memory packets of a request go to \col{NoC Mem Hops}, \col{NoC Mem Wait} and \col{NoC Mem Ahead} instead, and log as \code{MEM\_BUS}.};
\end{tikzpicture}
\end{adjustbox}
\caption{What one packet leaves behind: a bus record from injection to ejection, one record per
router it crosses, and the trackers added to its request's report row. Illustrated with the GetS
of \cref{fig:noc-example} delayed four ticks at router 1 by an older write-back.}
\label{fig:noc-observe}
\end{figure}

\paragraph{The report.} The NI stamps \code{REQ\_BUS} or \code{RESP\_BUS} EXIT at ejection, where
the mesh stamps it. The \code{LatencyReport} bus columns therefore span a packet's whole traversal,
and the event-path tiling of \cref{ch:request-end-to-end} holds. For a supply that goes to the
requester and to the home, the requester's copy sets the stamp. With a \code{Network} the report
gains five columns, accumulated over the request's coherence packets (\cref{fig:noc-observe}):
\col{NoC Hops}, \col{NoC Wait}, \col{NoC Transfer}, \col{NoC Ahead} and \col{NoC Ahead Winner}.
\col{NoC Transfer} is what the packets would take unloaded, \col{NoC Wait} the rest, and
their sum is the packets' injection-to-ejection time. \code{Summary.csv} gains \col{Worst-case NoC
Wait} and \col{Worst-case NoC Ahead}. With \key{memory\_on\_network} three more columns,
\col{NoC Mem Hops}, \col{NoC Mem Wait} and \col{NoC Mem Ahead}, collect the request's memory read
and fill. Reports of runs without a \code{Network} keep their columns.

\paragraph{The raw trace.} With the trace on (\cref{ch:monitoring}) every router writes records
under the Logger role \code{NOC} with \code{comp} = router id: ENTER when a head enters an input VC
(flags: input port and VC), SERVICE when it wins VC allocation (output port and VC), EXIT when the
tail leaves (output port and flit count). The bus records of a packet carry the destination id in
their flags, which tells the copies of a multicast apart.

\paragraph{The viewer.} On a network run the visualizer draws one lane per router and one array
lane per LLC bank, outlines the selected request's router stays and joins them into its route, one
chain per packet copy. A \emph{Network contention} grid under the lanes shows, per router and input
port, the cycles packets stayed beyond the shortest stay of their kind there, which is time spent
waiting for a VC, the switch or a credit.

\paragraph{\code{NetworkStats.csv}.} At the end of a run, normal or not, the network writes
\file{NetworkStats.csv} next to the reports, with the columns \code{section}, three keys and four
values (\cref{tab:noc-stats}).

\begin{table}[htbp]\centering\small
\caption{Sections of \code{NetworkStats.csv}.}\label{tab:noc-stats}
\begin{tabularx}{\textwidth}{@{}lY@{}}
\toprule section & keys and values \\ \midrule
\code{port} & router, port, port name: flits sent, ticks busy, switch-allocation attempts, losses \\
\code{class} & class: packets, mean and maximum injection-to-ejection latency, maximum wait \\
\code{classtime} & class: first injection tick, last ejection tick \\
\code{endpoint} & endpoint: messages delivered to its RX, popped by it, pushes refused, packets received \\
\code{received} & endpoint: packets, mean and maximum latency, owner of the first packet \\
\code{ahead} & endpoint: the largest \col{NoC Ahead} among its packets and the owner that won most \\
\code{epclass} & endpoint, class: packets received in that class \\
\code{generator} & generator id: sent, received, retries, out-of-order deliveries \\
\code{totals} & flits and packets injected and ejected and still in flight, flit-link traversals, ticks \\
\code{config} & topology, routing (as resolved from \code{auto}), \code{flit\_bytes}, \code{vc\_depth}, \code{router\_latency}, \code{link\_latency} \\
\code{options} & \code{vcs\_per\_class}, \code{escape\_vc}, \code{switch\_latency}, \code{sa\_allocator}, \code{routing\_seed}, \code{arbiter\_type} \\
\code{routerlatency} & router id: its \key{router\_latency[<router id>]} override, one row per override \\
\bottomrule
\end{tabularx}
\end{table}

\section{Errors}\label{sec:noc-errors}

Every network error prints \code{Network:} followed by what went wrong, with the router, port,
packet and message id where they apply, runs the hang dump when it is enabled, and exits with
status 2. At start-up the network refuses, with the key named: an invalid value of any key in
\cref{tab:noc-keys}, an endpoint without a router, an asymmetric \code{custom} matrix, an unknown
topology name (the message lists the registered builders) or routing name, a \code{routing\_table}
entry naming a missing port, a \code{routing\_via} or \code{link\_weight} entry naming a router
that is not a neighbour, a link weight below 1, a \key{router\_latency[<router id>](i)} for an id that is
not a router, a router-numbered key whose brackets do not hold an integer, a routing algorithm on a topology or VC count it does not support
(\cref{tab:noc-routing}), an \key{escape\_vc} not below \key{vcs\_per\_class}, an ordering
requirement the routing cannot meet, a cyclic channel dependency graph or escape network, the
deferred real-time keys, and a protocol that fails the check. At run time these are fatal: a destination that is not an endpoint,
an unclassifiable message, a routing result naming a missing port or an empty VC mask, a VC
overflow, a credit counter above \key{vc\_depth}, and a delivery without an RX reservation. None of
them can happen in a correct configuration. They exist so that an error in a new routing algorithm
or protocol shows itself at once.

\section{Configuration keys}\label{sec:noc-keys}

\Cref{tab:noc-keys} lists the network keys with their defaults, under \code{interconnect.} in a
system file. \Cref{tab:noc-syskeys} lists the system and controller keys the network features use.
\Cref{app:config-keys} has every shipped file in full.

\begin{longtable}{@{}p{0.29\textwidth}p{0.16\textwidth}p{0.49\textwidth}@{}}
\caption{Network keys (\code{configuration/Interconnect/Network.csv}).}\label{tab:noc-keys}\\
\toprule key & default & meaning \\ \midrule\endfirsthead
\toprule key & default & meaning \\ \midrule\endhead
\bottomrule\endlastfoot
\multicolumn{3}{@{}l}{\textit{topology}}\\
\key{topology(s)} & \code{custom} & \code{custom}, \code{grid2d}, \code{torus2d}, \code{ring}, \code{star}, \code{mesh3d} or \code{torus3d} \\
\key{component\_ids(vi)} & from \code{Mesh.csv} & the endpoint ids \\
\key{dims(vi)} & 0 & $k_x,k_y$ for \code{grid2d} and \code{torus2d}, $k_x,k_y,k_z$ for \code{mesh3d} and \code{torus3d}, $k$ for a ring \\
\key{node\_map[<id>](i)} & none & the router of an endpoint, on generated topologies \\
\key{connection\_matrix[i](vi)} & from \code{Mesh.csv} & \code{custom}: node id, then its neighbours \\
\key{router\_ids(vi)} & $-1$ & \code{custom}: the explicit routers. $-1$: every node that is not an endpoint \\
\key{local\_ports(i)} & 0 & local ports per router. 0: one per endpoint \\
\key{home\_ids(vi)} & $-1$ & the homes, for the classifier. $-1$: set by the system from the caches' home maps \\
\multicolumn{3}{@{}l}{\textit{packets, buffers and timing}}\\
\key{flit\_bytes(i)} & 16 & link width in bytes. 0: one flit per packet \\
\key{header\_bytes(i)} & 8 & bytes charged for a packet header \\
\key{num\_classes(i)} & 3 & 3: request, forward, response. 2: the legacy split by data \\
\key{vcs\_per\_class(i)} & 1 & VCs per class per input port \\
\key{vc\_depth(i)} & 4 & flits per VC \\
\key{router\_latency(i)} & 2 & ticks from a head's arrival to its earliest VC allocation, at least 0 \\
\key{router\_latency[<router id>](i)} & none & overrides \key{router\_latency} at one router \\
\key{switch\_latency(i)} & 0 & ticks from a head's VC grant to its earliest switch allocation, at least 0 \\
\key{link\_latency(i)} & 1 & ticks per link, at least 1. Local links cost 0 \\
\key{class\_link\_latency(vi)} & $-1$ & one link latency per class. $-1$: \code{link\_latency} \\
\key{credit\_latency(i)} & $-1$ & ticks for a credit to return. $-1$: \code{link\_latency} \\
\key{link\_model(s)} & \code{pipelined} & \code{pipelined}: one flit per tick per link. \code{serial}: a flit holds the link for its latency \\
\key{half\_duplex(i)} & 0 & 1: the two directions of a connection share one link \\
\key{inject\_queue\_depth(i)} & 16 & packets per class at the NI. At least the number of caches \\
\key{eject\_queue\_depth(i)} & 16 & messages per class in the endpoint's RX \\
\key{rx\_serving(s)} & \code{priority} & \code{priority}: response, forward, request. \code{alternate}: the legacy rotation \\
\key{m\_clk\_period(i)} & 50 & from \code{Mesh.csv}. One tick per two periods, one core cycle \\
\multicolumn{3}{@{}l}{\textit{routing and allocation}}\\
\key{routing(s)} & \code{auto} & \code{auto} (\code{XY} on grids, tori, 3-D lattices and on a ring with 2 VCs or more, else \code{Table}), or an algorithm of \cref{tab:noc-routing} \\
\key{routing\_via[<router>](vi)} & none & \code{Table}: \code{dest,neighbour} pairs overriding the computed table \\
\key{link\_weight[<router>](vi)} & none & \code{Table}: \code{neighbour,weight} pairs. Any entry switches to weighted shortest paths \\
\key{routing\_table[<router>](vi)} & none & \code{Table}: \code{dest,port} pairs, applied last \\
\key{routing\_seed(i)} & 1 & \code{O1TURN}, \code{Valiant}, \code{ROMM}: the seed of the per-packet draw \\
\key{escape\_vc(i)} & $-1$ & \code{AdaptiveMinimal}: the escape VC, and on wrapped lattices also the next one. $-1$: none \\
\key{arbiter\_type(s)} & \code{FCFSArbiter} & injection, VC and switch allocation: \code{FCFSArbiter} or \code{RRArbiter} \\
\key{sa\_allocator(s)} & \code{output\_first} & switch allocation: \code{output\_first} or \code{input\_first} \\
\key{deadlock\_check(s)} & \code{refuse} & the protocol-and-network check: \code{refuse} or \code{warn} \\
\end{longtable}

\begin{table}[htbp]\centering\small
\caption{System and controller keys of the network features.}\label{tab:noc-syskeys}
\begin{tabularx}{\textwidth}{@{}llY@{}}
\toprule key & where & meaning \\ \midrule
\key{interconnect\_type(s)} & system & \code{Network} selects this chapter's network in \code{MultiCoreSystem\_Mesh} \\
\key{requires\_ordering(s)} & system & \code{pair\_per\_class} or \code{none}: what the protocol needs from the routing \\
\key{llc\_banks(i)} & system & number of LLC banks, configured as \code{llc\_controller[k]} when above 1 \\
\key{memory\_on\_network(i)} & system & 1: memory controllers and the banks' memory ports are endpoints \\
\key{main\_memory\_type(s)} & system & \code{MainMemoryController} (default) or \code{MCsim}, one per memory controller \\
\key{mcsim\_scheduler(s)} & system & the MCsim scheduler, default \code{FRFCFS} \\
\key{home\_ids(vi)} & cache & the LLC banks. Absent or $-1$: \code{m\_shared\_memory\_id} \\
\key{home\_interleave\_shift(i)} & cache & line bits skipped by the home map, default 0 \\
\key{mem\_port\_id(i)} & LLC bank & the bank's memory-side endpoint \\
\key{mem\_ids(vi)} & LLC bank & the memory controllers. Absent or $-1$: \code{m\_shared\_memory\_id} \\
\key{mem\_interleave\_shift(i)} & LLC bank & line bits skipped by the memory map, default 0 \\
\bottomrule
\end{tabularx}
\end{table}

\section{Validation}\label{sec:noc-validation}

One runner builds the simulator and runs the network tests:

\begin{lstlisting}[style=shell]
bash test/network/run.sh all     # every test
bash test/network/run.sh 4       # one step: 0..8, M2..M6, D, T
bash test/network/run.sh reg     # the regression suite only
\end{lstlisting}

Each test prints \code{PASS}, \code{FAIL} or \code{XFAIL} with its id and writes a row to
\file{test/network/out/results.csv}. A single step runs the regression suite first. The timing
tests take their expected values from the formulas of this chapter, not from earlier runs, so a
wrong engine cannot confirm itself. REG-MESH, REG-TDM, REG-MEMNET, MEM-OFF and S6 compare against
committed references instead. The trace checks need a Python with numpy and duckdb. S4-ASAN,
S5-ASAN, S5-XPLAT and S8-BOOKSIM run only where \code{wsl.exe} and the sanitizer build
\file{build_asan} or BookSim exist, and report XFAIL elsewhere. \Cref{tab:noc-tests} summarises what the suite establishes.

\begin{longtable}{@{}p{0.24\textwidth}p{0.71\textwidth}@{}}
\caption{What the test suite (\code{test/network/}) checks.}\label{tab:noc-tests}\\
\toprule tests & what they establish \\ \midrule\endfirsthead
\toprule tests & what they establish \\ \midrule\endhead
\bottomrule\endlastfoot
S1, S2, S3 (unit) & classifier, flit counts, every topology generator and port numbering, every routing algorithm, the channel dependency check (a ring under \code{Table} is refused), the ordering refusal, and both arbiters, whose network and bus entry points agree on 1000 random elections \\
S4 unloaded & corner to corner on 2$\times$2 and 4$\times$4 grids, a torus and a star, for 1 and 5 flits: 8 and 12, 20 and 24, 8 and 12, 2 and 6 ticks, each equal to the formula \\
S4 flow control & per-VC throughput 0.5, 1 and 1 flit per tick for \code{vc\_depth} 1, 2, 4, and 0.5 with \code{credit\_latency} 3. Wormhole blocking delays an unrelated packet with shallow VCs and not with deep ones. Half duplex halves a connection's throughput (404 against 204 ticks) \\
S4 conservation & every run ends with flits injected equal to flits ejected, no credit above its initial value, and three reruns byte-identical \\
S4 saturation & 4$\times$4, \code{XY}, 5-flit packets, uniform traffic saturates at an offered 110 packets per 1000 ticks per endpoint. Hotspot (20) and transpose (80) saturate earlier \\
S4 ordering, faults & no reordering under \code{XY} and \code{Table}. Adaptive routing reorders (an expected failure that reports the count). Each fatal path exits with status 2 \\
S5 system & directory MSI on \code{Grid}: all ten EEMBC traces complete, event-path tiling with no failure, byte-identical reruns, the placement figures of \cref{sec:noc-example}, and the hang dump on a deliberately wedged ring. Where WSL and the sanitizer build exist: identical Windows and Linux runs and a clean AddressSanitizer run \\
S6 legacy & \code{MultiCoreSystem\_NetNoC} against the legacy \code{NoC} references: S6-IDENTITY (byte identity), S6-FALLBACK (the per-core \code{Summary.csv} deltas) and S6-NOC-ALIAS, all three XFAIL by the accepted-divergence file \file{test/network/ref/s6-accepted.txt}. The suite asserts no per-path figures. Those of \cref{fig:noc-legacy} come from \file{docs/NoC_Implementation_Report.md} \\
S7 observability & router hops in the trace equal \col{NoC Hops}, \col{NoC Wait} plus \col{NoC Transfer} equals the network time of every request, and the contention example reports a wait of 4 and the right winner \\
S8 BookSim & 4$\times$4 mesh, dimension-order routing, one VC, 4-flit buffers, 5-flit uniform traffic, with BookSim's buffer residency mapped onto \code{router\_latency} 3 and \code{credit\_latency} 4. The test passes when the saturation throughputs agree within 10\,\%. The recorded measurement is 3.5\,\% (0.288 and 0.279 flits per cycle per node) \\
P2 protocol & three classes with an RX depth of 2 under both arbiters, \code{inject\_queue\_depth} 4 without a refused push (2 wedges and exits with status 2 as documented), ADMIT-RF on three controllers, the protocol check passing on three classes and refusing two, and every request reaching the bank the home map names \\
MEM & memory on the network: existing presets unchanged when it is off, EEMBC complete with one and two controllers and with MCsim, no read-after-write overtake, every memory packet in its class, memory hops equal to the formula, and every read and write reaching the controller the map names \\
D & directory MESI and MOESI on every grid preset, under both arbiters and with small queues: every run completes without a fault \\
T (unit) & topology plug-ins: the port order of every builder and \code{custom} row order, \code{routing\_via}, the builder registry, lattice ports, the dependency check over the candidates of every algorithm on several grids and VC counts, the escape-network check (the one-escape-VC torus is refused), the \code{AdaptiveMinimal} fallback, 3-D lattices, \code{XY} on rings of 3 to 9, weighted \code{Table}, and the oblivious draws (minimal for \code{O1TURN} and \code{ROMM}, varying with the packet and the seed) \\
T system & T4-UNLOADED: a ring of 8, \code{mesh3d} 4$\times$2$\times$2 and \code{torus3d} 3$\times$3$\times$3, 1 and 5 flits, each equal to the formula. T4-AUTO-RING: \code{auto} gives \code{XY} on a ring with 2 VCs, \code{Table} on a ring of 3, and a ring of 8 with one VC is refused. T4-ROUTER-OPTIONS: \code{switch\_latency} 1 (27 ticks), \code{router\_latency[3]} 6 (24) and \code{input\_first} (24) on a 4$\times$4 grid, each equal to the formula. T4-SOAK, T4-SOAK-2: loaded runs of every new topology and algorithm, \code{AdaptiveMinimal} on grid, torus, ring and 3-D mesh, \code{input\_first} with \code{RRArbiter}, the new patterns and on/off injection complete with conservation. T4-PATTERNS: under \code{bitcomp}, \code{bitrev}, \code{shuffle} and \code{tornado} every endpoint hears from exactly its partner. TF-RANDPERM-NO-FIXED-POINT: under \code{randperm} every endpoint sends and receives. TF-REFUSALS: a malformed router id in a key and a bit pattern on 12 endpoints end with status 2 and the reason. TF-STATS-OPTIONS: \code{NetworkStats.csv} has the \code{options} and \code{routerlatency} rows. T5-ADAPT-TORUS-REFUSED \\
REG & REG-BUS, REG-MESH, REG-TDM and REG-MEMNET: the bus, legacy mesh, TDM and memory-on-network references stay byte-identical. REG-ENV: the environment check's sample run gives its known request count and worst DRAM latency. REG-HAZARD: no memory read overtakes an older write of its line \\
\end{longtable}

The saturation point of the shipped defaults is higher than BookSim's (0.48 against 0.28 flits per
cycle per node with the same buffers), because an Octopus slot is freed the tick its flit is
forwarded. That is the faster router the defaults describe. To compare with a router that holds a
slot through switch allocation and traversal, raise \key{router\_latency} and
\key{credit\_latency} as above.

\section{Limitations}\label{sec:noc-limits}

\begin{itemize}
\item Only the directory protocols run on the network. Snooping needs an ordered broadcast, which
  the network does not provide.
\item No real-time arbitration: \code{TDMArbiter}, \code{PriorityArbiter}, \code{TDM\_XY},
  \code{slot\_table} and \code{vc\_priority} are refused as deferred.
\item A ring of four or more routers needs \code{XY} with \key{vcs\_per\_class} of at least 2.
  With one VC, \code{auto} falls back to \code{Table}, which has no dateline and is refused there.
\item The turn models run on \code{grid2d} only, \code{O1TURN} on 2-D lattices without wraps, and
  \code{Valiant} and \code{ROMM} on lattices without wraps. No oblivious algorithm runs on a torus
  or a ring.
\item On a torus or a ring \code{XY} uses VCs 0 and 1 only. More VCs per class add buffers that
  \code{XY} leaves empty.
\item The \code{transpose} pattern needs a 2-D lattice. Elsewhere every generator sends to the
  first other endpoint. The bit patterns need a power-of-two number of endpoints.
\item On a network of more than 64 routers the start-up check trusts the algorithms that are
  deadlock-free by construction and does not build their dependency graph. \code{Table} and the
  escape network of \code{AdaptiveMinimal} are checked at every size.
\item Every link has a reverse link, so a butterfly cannot be expressed, and an endpoint attaches
  to exactly one router.
\item Injection follows age only among eligible packets. Under \code{FCFSArbiter} a younger packet
  of another class leaves first when the older one finds no free VC with a credit. Under
  \code{RRArbiter} the local port rotates over packet owners, so an endpoint's packets of different
  classes may leave out of order. The directory FSMs tolerate the cross-class reordering the fabric
  causes anyway.
\item Under \code{RRArbiter} a candidate whose owner is not an endpoint id never wins an election
  unless it is alone. No shipped preset produces such an owner.
\item A local port carries one flit per tick in each direction, and \key{header\_bytes} is the same
  for every class.
\item Multicast is one packet per destination. There is no multicast tree.
\item The NI cannot take back a message once pushed (\code{rollback} finds nothing).
\item The legacy mesh keys \key{buffers\_max\_size}, \key{m\_request\_latency},
  \key{m\_data\_latency} and \key{switch\_ids} are not read by the network. Use the queue depths, \key{class\_link\_latency}
  and \key{router\_ids}.
\item A read filled from a pending write costs no cycles, at the bank and at the controller.
\item There is no power or area model.
\end{itemize}

\begin{wherebox}
\begin{itemize}[nosep,leftmargin=1.2em]
\item \file{header/Interconnect/Network/}, \file{src/Interconnect/Network/}: \code{Network} (the engine, statistics, hang dump), \code{NetworkInterface}, \code{Topology} and \code{TopologyBuilder} (the graph, the lattice capability, the builder registry), \code{Routing} (the algorithms, the channel dependency check), \code{Classifier}, \code{ProtocolCheck}, \code{TrafficGenerator}
\item \file{src/Interconnect/Network/topologies/}: the builders (\file{Lattice.cpp}, \file{Graph.cpp}) and their registry (\file{Builtin.cpp})
\item \file{header/Arbiters/Arbiter.h}: \code{Candidate} and \code{electIndex}
\item \file{src/SystemConfigurations/MultiCoreSystem_Mesh.cpp}: wiring, LLC banks, memory on the network. \file{src/SystemConfigurations/NetworkOnly.cpp}: network-only mode
\item \file{src/CacheControllers/CacheController.cpp}: transmit-side admission (\code{txNeed}, \code{txAdmit})
\item \file{configuration/Interconnect/Network.csv}, \file{configuration/TrafficGenerator.csv}, the presets in \file{configuration/SystemConfigurations/}
\item \file{test/network/}: the runner, the unit tests, the checks and the BookSim configuration
\item \file{docs/NoC_Architecture_Spec.md}, \file{docs/NoC_Implementation_Report.md}: the specification and the measured record of each step. \file{docs/NoC_Topology_Plugin_Plan.md}: the plug-in design
\end{itemize}
\end{wherebox}
\endgroup
